\documentclass[11pt]{amsart}

\input{preamble}

\addbibresource{references.bib}

%%% Project-specific macros %%%
\nc{\A}{\mathbb{A}}
\nc{\OO}{\mathcal{O}}
\nc{\Gr}{\mathrm{Gr}}
\nc{\Hilb}{\mathrm{Hilb}}
\nc{\IC}{\mathrm{IC}}
\nc{\Av}{\mathrm{Av}}
\nc{\bv}{\mathbf{v}}
\nc{\bw}{\mathbf{w}}
\nc{\bb}{\mathbf{b}}
\nc{\one}{\mathbf{1}}
\nc{\pre}[1]{\cite[#1]{openaiEquivariantCohomological2026}}
\nc{\Pol}{S}
\nc{\Zm}{Z_m}

% Details: routine verifications that a reader could reconstruct.  They are
% hidden by default; change \detailsfalse to \detailstrue to show them.
\usepackage{comment}
\newif\ifdetails
\detailsfalse
\ifdetails
\newenvironment{details}{\par\color{black!60!green}\noindent\textbf{Details.}\ }{\par}
\else
\excludecomment{details}
\fi

\title[Hikita's conjecture for quiver gauge theories]{Hikita's conjecture for quiver gauge theories:\\ after the slop}


\begin{document}

\begin{abstract}
Let $X$ be a Nakajima quiver variety, with loops and multiple arrows
allowed, at a generic stability parameter, and let $A$ be the Coulomb
branch of the same gauge theory.  Using the comparison map announced by
Botta and Tamagni, we explain a proof that the B-algebra of $A$ for the
cocharacter given by the stability parameter is isomorphic to $H^*(X)$,
equivariantly for flavor symmetries, and that the B-algebra of the quantized
Coulomb branch is isomorphic to $H^*_{\C^\times}(X)$.  The theorem and the
overall strategy are due to \cite{openaiEquivariantCohomological2026}.  The main change is a
new proof of the dimension count for quivers with a loop at every vertex, by
restriction to Levi subgroups, which does not use Hausel's formula.  On the
quiver variety side this count becomes the decomposition theorem for the map
recording the characteristic polynomials of the loops.
\end{abstract}

\maketitle

\section{Introduction}\label{sec:intro}

These notes are aimed at producing a human readable and checked proof of the Hikita conjecture for quiver gauge theories.  This builds on the claimed proof in the ``slop dump'' of Oct.\ 7, 2026
 \cite{openaiEquivariantCohomological2026}, an unrefereed
preprint by OpenAI dated 24 September 2026.  

I should emphasize that this article is {\it not} a simple rewriting of the argument in that document, nor does it constitute an endorsement of the lazy approach of that dump, or the many profound violations of research ethics involved in it.  It was a deeply irresponsible and selfish action carried out without regard to the consequences for individuals, or the very research apparatus that, of course, was a necessary precondition for the existence of the AI system that produced it.  

However, given the existence of a claimed proof, I think it is in the interest of the mathematical community as a whole to have access to a version of the argument that is readable and checkable by humans.  Given the difficulty of understanding the original argument of the OpenAI preprint, the only practical way of doing this was to use AI itself, but in this case guided by a human who understands the field.  Thus, this paper was written with considerable assistance from Claude in understanding the arguments of the OpenAI preprint, finding references in the literature, and rewriting the argument in a more readable form.   

\Ben{At the moment I don't want to say that humans have verified the full argument, since that isn't true.  I *will*.}

Fix a finite quiver $Q$, loops and multiple arrows allowed, and dimension
and framing vectors $\bv,\bw\in\Z_{\ge0}^{Q_0}$.  Put
\[
 G=\prod_{i\in Q_0}\GL(V_i),\qquad
 N=\bigoplus_{a\colon i\to j}\Hom(V_i,V_j)\oplus\bigoplus_i\Hom(W_i,V_i).
\]
Let $F$ be a torus acting on $N$ and commuting with $G$, and let
$\C^\times_\hbar$ act on $T^*N$ with weight $\frac12$ on both $N$ and $N^*$
(through its double cover), so that the symplectic form has weight $1$. 

For a choice of character $\nu\in\Z^{Q_0}$, let $X_\nu$ be the Nakajima quiver variety with this stability parameter, and let $A_{F,\hbar}$ be the quantized Coulomb branch of $(G,N)$ with flavor symmetry $F$.  Let $B^\nu(A_{F,\hbar})$ be the B-algebra of $A_{F,\hbar}$ for the action induced by the cocharacter $\nu$ of $\pi_1(G)^\vee$.
\begin{theorem}\label{thm:main}
Let $\nu$ be generic.  Then we have an isomorphism of graded algebras
\[B^\nu(A_{F,\hbar})\cong H^*_{F\times\C^\times_\hbar}(X_\nu).\]
\end{theorem}
The existence of a surjective map $B^\nu(A_{F,\hbar})\to H^*_{F\times\C^\times_\hbar}(X_\nu)$ follows immediately from work of Botta and Tamagni \cite{bottaQuasimapCritical2026}. One of the outrageous aspects of the OpenAI preprint is that they effectively rewrote the argument of Botta and Tamagni's paper without giving them proper credit (and that OpenAI clearly did not apply sufficient due diligence to avoid this happening in many of their preprints).  
We only briefly summarize this construction.

Applying the graded Nakayama lemma, to complete the proof of Theorem~\ref{thm:main}, we only need to show an inequality 
\[\dim B^\nu(A)\leq \dim H^*(X_\nu).\]

Together with the localization theorem, this lemma gives us a useful reduction: if \cref{thm:main} applies for a given choice of $Q$, $\bv$, $\bw$, $F$ and $\nu$, then every choice of point in the Lie algebra $\ff$ of $F$ gives a vector field on $X_\nu$ whose fixed point set is a union of quiver varieties for gradings of the same quiver, and \cref{thm:main} applies to each of these fixed point components.  

This brings us to one of the genuinely interesting approaches taken in the OpenAI preprint: they reduce the proof of the dimension inequality to the case of quivers with a loop at every vertex; we have consulted with many experts, and none can point to a place in the literature where this approach was used.  I wish I understood the importance of this reduction better, but we find two explicit results which might help the reader see how it is helpful.  In this case we have an explicit formula for a nonnegative integer $H(Q,\bv,\bw)$ which we need to verify
\begin{enumerate}
\item $\dim B^\nu(A)\le H(Q,\bv,\bw)$.  This comes from a filtration of
$B^\nu(A)$ whose graded pieces are sums, over multipartitions of $\bv$, of
quotients of the invariants of finite groups acting on the cohomology rings
of hypertoric varieties (Proposition~\ref{prop:filt}, which is
\pre{Prop.~5.2}).  Here the loops are used on the Coulomb side: they supply
the roots of $G$ as weights of $N$, so that the Coulomb branch is the
Weyl-invariant part of an abelian Coulomb branch, with no root denominators
(Proposition~\ref{prop:teleman}).
\item $H(Q,\bv,\bw)\le\dim H^*(X_\nu)$.  The preprint deduces this from
Hausel's formula \cite{hauselBettiNumbers2006} for the Betti numbers of $X_\nu$,
together with an identity for the cohomology of toric varieties proved by
Kirwan's method \pre{\S6}; Remark~\ref{rem:hausel} explains how these fit
together, following notes of Nakajima \cite{nakajimaMultipartitionSpaces2026}.  We
give a different argument, which does not use Hausel's formula.  

The filtration of $B^{\nu}(A)$ by subspaces $F_L$, indexed by the number $L$ of parts of a multipartition, induces a filtration of $H^*(X_\nu)$ by subspaces $F'_L$.  Thus, we only need to show that the induced map $F_L/F_{L-1}\to F'_L/F'_{L-1}$ is injective.  We prove this by reduction to a Levi subgroup $M$ of $G$; this gives us a new quiver variety $X(M)$, and turns the piece of $F_L/F_{L-1}$ for a multipartition $m$ into a level module over $H^*_M(\mathrm{pt})$.  Of course, we only need to show injectivity on its socle, which is its top degree piece.  Thus, we only need to show that $H^*(X(M))$ has a sufficiently large top degree piece, and ultimately we construct by hand enough distinct core components of $X(M)$ to show this.
\end{enumerate}
Here
\[
 H(Q,\bv,\bw)=\sum_m\dim\big(\Pol_m/J_m\big)^{\cW_m},
\]
the sum over multipartitions $m$ of $\bv$, where $\Pol_m/J_m$ is the
cohomology ring of a hypertoric variety and $\cW_m$ is a product of symmetric
groups acting on it (\S\ref{ss:notation}).  For example, for the Jordan quiver
with $\bw=1$, so that $X_\nu=\Hilb^n(\C^2)$, every summand is $1$, and
$H(Q,n,1)$ is the number of partitions of $n$.  Both sides of~(1) and~(2) are
independent of the generic parameter $\nu$ (Lemmas~\ref{lem:generic}
and~\ref{lem:betti}), so we work at determinant stability.  Keeping track of
degrees gives the Poincar\'e polynomial of $X_\nu$ as well
(Theorem~\ref{thm:loops}); Example~\ref{ex:A2loops} works out a small case.

The loops can then be removed.  Given $Q$, add a loop $\ell^+_k$ at every
vertex, with a flavor torus $\C^\times_u$ scaling it.  As above, the quiver
variety $X_\nu$ of $Q$ is a component of the $\C^\times_u$-fixed locus of the
quiver variety of the new quiver, the one where the added loops vanish.  The
localization theorem, and a matching statement on the Coulomb side after
localizing at $u=1$, transfer the theorem from the new quiver to $Q$
(Proposition~\ref{prop:loopremoval}, which is \pre{Prop.~7.3}).

On the quiver variety side, the multipartitions have a direct meaning.  The
characteristic polynomials of the loops give a map $\chi$ from $X_\nu$ to a
space of tuples of polynomials, and the summand of $m$ in $H(Q,\bv,\bw)$ comes
from the stratum where the eigenvalues of the loops have multiplicities given
by $m$: the decomposition theorem for $\chi$, read at one point, gives
$H^*(X_\nu)$ as a sum over multipartitions (Theorem~\ref{thm:higgs}).  We do
not use this in the proof.


\subsection{Set-up}\label{ss:setup}
We now make the objects above precise.  With $\C^\times_\hbar$ acting on $T^*N$
as above, on the Coulomb side $\hbar$ is the equivariant parameter of loop rotation
combined with the scaling of $N$ with weight $\frac12$, as in
\cite[\S\S2.2, 5.7]{hilburnBFNSpringer2023}.  Other conical actions in the literature, such as
that of \cite[\S7.2]{KTWWY}, differ from this one by a rescaling and a
cocharacter of the torus scaling the framing, which shifts the flavor
parameters by multiples of $\hbar$.  Put $R=H^*_{G\times
F}(\mathrm{pt})$ and $R_\hbar=R[\hbar]$, and let $\mu\colon T^*N\to\mathfrak g^*$
be the moment map.

\emph{Stability.}  A point of $T^*N$ is a representation of the doubled
quiver of $Q$ on the $V_i$, together with maps $i_k\colon W_k\to V_k$ and
$j_k\colon V_k\to W_k$ at each vertex $k$.  Following \cite[\S1]{crawley-boeveyGeometryMoment2001}, we regard it as a
representation of dimension $(\bv,1)$ of the doubled quiver of $\hat Q$,
which is $Q$ with one more vertex $\infty$ and $w_i$ arrows from $\infty$ to
$i$.  For $\nu\in\Z^{Q_0}$ put $\hat\theta=(\nu,-\nu\cdot\bv)$.  A point is
\emph{$\nu$-semistable} if $\hat\theta\cdot\dim S\ge0$ for every
subrepresentation $S$, and \emph{$\nu$-stable} if moreover equality holds
only for $S=0$ and for the whole representation \cite{kingModuliRepresentations1994}.  The same
inequalities define $\theta$-stability, and the quiver variety $X_\theta$
below, for real $\theta\in\mathbb R^{Q_0}$.  Call $\nu$ \emph{generic} if
$\nu\cdot\bb\neq0$ for all $0<\bb\le\bv$.  Then every $\nu$-semistable point is
stable, with trivial stabilizer in $G$, and the semistable locus depends only
on the signs of the $\nu\cdot\bb$.
\begin{details}
The subrepresentations have dimensions $(\bb,0)$ and $(\bb,1)$ with
$0\le\bb\le\bv$, on which $\hat\theta$ takes the values $\nu\cdot\bb$ and
$-\nu\cdot(\bv-\bb)$.  So for generic $\nu$, $\hat\theta\cdot\dim S\neq0$ unless
$S$ is $0$ or the whole representation, and the automorphisms of a stable
point are scalars that are the identity at $\infty$.
\end{details}
For generic $\nu$ the quiver variety
$X_\nu=\mu^{-1}(0)^{\nu\text{-ss}}/G$ is smooth, and
$H^*_{G\times F}(\mu^{-1}(0)^{\nu\text{-ss}})=H^*_F(X_\nu)$.  The
tautological bundle $\cV_i$ on $X_\nu$ comes from the $G$-module $V_i$.  The
\emph{Kirwan map} $\kappa\colon R\to H^*_F(X_\nu)$ is restriction from
$H^*_{G\times F}(\mathrm{pt})$ to $H^*_{G\times F}(\mu^{-1}(0)^{\nu\text{-ss}})$; it
sends a symmetric polynomial in the Chern roots of the $V_i$ to the same
polynomial in the $F$-equivariant Chern roots of the $\cV_i$.  Determinant
stability $\nu_{\det}=(1,\dots,1)$ is generic.

This is stability for the character $\chi_\nu=\prod_i\det_{V_i}^{\nu_i}$, with
semi-invariants $f(g\cdot x)=\chi_\nu(g)^nf(x)$, as in
\cite[\S2.1]{hilburnBFNSpringer2023}; for $\nu_{\det}$ it says that the images of the $W_k$
generate $V$, by~\ref{I:cyc}.  The stability condition of Nakajima, under which no
nonzero invariant subspace lies in the kernel of the maps $j_k$, is
$\nu=(-1,\dots,-1)$ in this convention \cite[Lemma~3.8]{nakajimaQuiverVarieties1998}; compare
\cite[Rem.~7.12]{hilburnBFNSpringer2023}.

\emph{Coulomb branches.}  Let $A_{F,\hbar}$ be the quantized Coulomb branch
of $(G,N)$ with flavor symmetry $F$ \cite{BFN}, and $A_F=A_{F,\hbar}/\hbar$,
which is commutative.  Both are graded by $\pi_1(G)=\Z^{Q_0}$ through the
components of the affine Grassmannian; we call this degree the
\emph{charge}.  The charge-zero part contains $R_\hbar$, respectively $R$,
as the classes supported over the base point.  For a charge $c$ put
$\deg_\nu(c)=-\sum_i\nu_ic_i$, and let $A_{[k]}$ be the sum of the charge
components with $\deg_\nu=k$.  The \emph{B-algebra} is
\[
 B^\nu(A_{F,\hbar})=A_{[0]}\Big/\textstyle\sum_{k>0}A_{[-k]}A_{[k]},
\]
and $B^\nu_F$ is the same construction for $A_F$.  For $\hbar\neq0$ the
order of the factors matters: reversing it amounts to replacing $\nu$ by
$-\nu$.  Our $\nu$-degree is the grading of \cite[\S5.1.6]{hilburnBFNSpringer2023} for the
character $\chi_\nu$, and $B^\nu$ is the B-algebra of \cite[\S5.7]{hilburnBFNSpringer2023}
and \cite[\S3.5]{KTWWY}.  As $A_F$ is
commutative, $B^\nu_F$ is the
coordinate ring of the fixed-point scheme of the cocharacter $\nu$ of the
torus $\pi_1(G)^\vee$, which acts on $\Spec A_F$ through the charge.  The
fixed-point scheme of the whole torus has coordinate ring
$B_F=A_F/\langle A_F^{(c)}:c\neq0\rangle$, where $A_F^{(c)}$ is the part of
charge $c$; it is a quotient of $B^\nu_F$.  We write $A_0$, $B^\nu_0$ and $B_0$
when $F$ is trivial.


\emph{The statement.}  Our proof uses Theorem~\ref{thm:map}, the
comparison map announced by Botta and Tamagni.  Setting $\hbar=0$ in
Theorem~\ref{thm:main} gives $B^\nu_F\cong H^*_F(X_\nu)$; more precisely, the
maps $R\to B^\nu_F$ and $\kappa$ are surjective with the same kernel
(Theorem~\ref{thm:classical}).  This classical statement is the generic case
of \pre{Thm.~1.1}, which also allows $\nu$ on a hyperplane $\nu\cdot\bb=0$
when every semistable point is still stable.  We do not treat that case.
With $F$ trivial it is Hikita's conjecture
\cite{hikitaAlgebroGeometricRealization2016}.  Theorem~\ref{thm:main} is its
equivariant and quantized form, due to Nakajima \cite[Conj.~8.9]{KTWWY}; with
the conventions of \S\ref{ss:setup}, it is
\cite[Conj.~5.26]{hilburnBFNSpringer2023} for quiver gauge theories.

\emph{Known cases.}  Hikita proved the case with $F$ trivial for Hilbert
schemes of points in the plane, type A Spaltenstein varieties and smooth
hypertoric varieties \cite{hikitaAlgebroGeometricRealization2016}.  For ADE quivers it is
\cite[Thm.~8.1]{KTWWY}, whose proof uses an isomorphism of Zhu that is
known in types A, D and $E_6$ and in some cases of types $E_7$ and $E_8$
\cite[Rem.~2.18, Cor.~2.17]{dumanskiKtheoreticHikita2025}.  With flavor, the classical statement is known for the
Gieseker variety, the Jordan quiver with framing $r$
\cite[Thm.~1.14]{krylovHikitaNakajimaConjecture2025}, and for type A quivers and the same ADE
cases \cite[Cors.~2.15, 2.17]{dumanskiKtheoreticHikita2025}.  The quantized statement is known for
smooth hypertoric varieties \cite[Thm.~1.2, Rem.~5.3]{kamnitzerQuantumHikita2021}.  Dinkins, Karpov
and Krylov prove a refined quantum Hikita conjecture for ADE quivers framed at
minuscule nodes, including all type A quivers and the Jordan quiver
\cite[Thm.~B]{dinkinsQuantumHikita2026}.

\subsection{The argument}\label{ss:route}
\S\ref{sec:map} shows that $R_\hbar\to B^\nu(A_{F,\hbar})$ is onto and states
the comparison map.  \S\ref{sec:loops} proves Theorem~\ref{thm:loops}, the case
of a loop at every vertex, at determinant stability and without flavor, by
establishing~(1) and~(2) above.  \S\ref{sec:deform} adds flavor and $\hbar$ by
graded Nakayama (Lemma~\ref{lem:flavor}), removes the loops
(Proposition~\ref{prop:loopremoval}) and changes the stability parameter
(Lemma~\ref{lem:betti}); this proves Theorem~\ref{thm:main}
(Theorem~\ref{thm:classical}).  \S\ref{sec:higgs} describes \S\ref{sec:loops}
on the quiver variety side: the multipartition filtration becomes the support
filtration of the decomposition theorem for $\chi$.  It is not used
elsewhere.  \S\ref{sec:bps} gives a second proof of
Proposition~\ref{prop:loopremoval}, which uses neither \S\ref{sec:loops} nor
added loops, but depends on results on cohomological Hall algebras whose
proofs are only sketched in the literature.

\emph{Status.}  The argument is complete given the cited results and
Theorem~\ref{thm:map}, which is announced in \cite{bottaQuasimapCritical2026}.
\S\ref{sec:status} lists what remains to be checked.

\subsection{Inputs}\label{ss:inputs}
We use the following results.
\begin{enumerate}[leftmargin=2.6em,label=(I\arabic*)]
\item\label{I:gen} \emph{Generation} \cite[Prop.~3.1,
Rem.~3.3]{weekesGeneratorsCoulomb2019}.  $A_{F,\hbar}$ is generated over $R_\hbar$ by the
dressed minuscule monopole operators $M_{\lambda,p}$, for dominant coweights
$\lambda$ of $G$ with all entries in $\{0,1\}$ or all in $\{0,-1\}$.  Here the
orbit $\Gr^\lambda$ is closed, $M_{\lambda,p}$ is the class of the part of
the BFN space of triples over $\Gr^\lambda$ multiplied by the dressing $p$, and
$p$ runs over $H^*_{G\times F}(\Gr^\lambda)=H^*_{T\times F}(\mathrm{pt})^{\cW_\lambda}$,
with $T\subseteq G$ the diagonal torus and $\cW_\lambda$ the stabilizer of
$\lambda$ in the Weyl group.  The same holds for $A_F$ over $R$.  The proof generates the associated graded of the
filtration by support in the orbit closures $\overline{\Gr^\Lambda}$
\cite[Prop.~6.8]{BFN}.  So the filtered piece $A_{\le\Lambda}$ is spanned
by words in the generators whose coweights sum to at most $\Lambda$.
\item\label{I:base} \emph{Base change} \cite[Prop.~3.22]{BFN}.
$A_{F,\hbar}$ is free over $H_{F,\hbar}=H^*_{F\times\C^\times_\hbar}(\mathrm{pt})$,
and $A_{F,\hbar}\otimes_{H_{F,\hbar}}\C=A_0$; likewise
$A_{F,\hbar}\otimes_{H_{F,\hbar}}H^*_F(\mathrm{pt})=A_F$.  The flavor parameters
and $\hbar$ have charge zero, and the B-algebra is a cokernel, so taking
B-algebras commutes with these base changes.
\item\label{I:kirwan} \emph{Kirwan surjectivity} \cite[Thm.~1.2,
Cor.~1.5]{mcgertyKirwanSurjectivity2018}.  $\kappa$ is surjective, $H^{\mathrm{odd}}(X_\nu)=0$, and
$H^*_F(X_\nu)$ is free over $H^*_F(\mathrm{pt})$.  The same holds with
$F\times\C^\times_\hbar$ in place of $F$.
\item\label{I:rot} \emph{Hyperk\"ahler rotation} \cite[Prop.~3.18]{bellamySymplecticResolutions2021}.
For generic $\theta\in\mathbb R^{Q_0}$, $X_\theta$ is homeomorphic to the
affine quotient $\mu^{-1}(-i\theta)/\!/G$.  Here $\C^{Q_0}$ is identified with
the center of $\mathfrak g^*\cong\mathfrak g$ by $\zeta\mapsto\sum_i\zeta_i\,\mathrm{id}_{V_i}$.  Loops are allowed.
\item\label{I:cyc} \emph{Determinant stability} \cite[Lemma~7.1]{hilburnBFNSpringer2023}.
A point of $T^*N$ is $\nu_{\det}$-semistable if and only if the images of the
$W_i$ generate $V=\bigoplus_iV_i$ under the arrows of the doubled quiver.
Loops are allowed.
\begin{details}
$\nu_{\det}\cdot\bb>0$ for $\bb\neq0$, so a subrepresentation of
dimension $(\bb,0)$ never destabilizes, and one of dimension $(\bb,1)$, which
contains the images of the $W_i$, destabilizes unless $\bb=\bv$.
\end{details}
\item\label{I:hypertoric} \emph{Hypertoric varieties}
\cite{hauselToricHyperKahler2002}.  Let a torus $S$ act linearly on a vector space $U$,
with weights $u_1,\dots,u_n$ (repeated according to multiplicity), none of them
zero.  Let $\mu\colon T^*U\to\mathrm{Lie}(S)^*$ be the moment map and $\theta$ a
generic character of $S$.  The \emph{hypertoric variety} is
$\mu^{-1}(0)/\!\!/_\theta S$.  A \emph{cocircuit} is a minimal set of weights
whose complement does not span the span of all the $u_i$.  Suppose the $u_i$
span $X^*(S)\otimes\Q$.  Then:
\begin{enumerate}[label=(\alph*)]
\item if the $u_i$ are unimodular, the hypertoric variety is smooth, of
dimension $2(n-\dim S)$ when it is nonempty \cite{hauselToricHyperKahler2002};
\item its cohomology ring is $H^*_S(\mathrm{pt})$ modulo the ideal generated by
the products $\prod_{u\in C}u$ over the cocircuits $C$
\cite[Thm.~1.1]{hauselToricHyperKahler2002};
\item this ring vanishes above degree $n-\dim S$, and it is level in the
sense of~\ref{I:alg}(c) \cite[Thm.~7.1, Prop.~6.7]{hauselToricHyperKahler2002},
\cite[Thm.~III.3.4]{stanleyCombinatoricsCommutative1983};
\begin{details}
The ring is $\C[\partial]/\mathrm{Ann}(\mathrm{vol}_1,\dots,\mathrm{vol}_r)$,
where $H^*_S(\mathrm{pt})=\C[\partial]$ acts by differentiation and the $\mathrm{vol}_i$ are
the volume polynomials of the maximal bounded regions of the hyperplane
arrangement of the $u_i$ and $\theta$, all homogeneous of degree $n-\dim S$
\cite[Thm.~7.1, Prop.~6.7]{hauselToricHyperKahler2002}.  An operator of degree above
$n-\dim S$ kills every $\mathrm{vol}_i$, and the ring is level by Stanley's theorem
\cite[Thm.~III.3.4]{stanleyCombinatoricsCommutative1983}, as in the proof of
\cite[Thm.~7.1]{hauselToricHyperKahler2002}.
\end{details}
\item the top Betti number is the top $h$-number of the matroid of the
arrangement \cite[Cor.~8.4]{hauselToricHyperKahler2002}, which is dual to the matroid of
the $u_i$; so it is the cographic matroid when the $u_i$ are the vectors of a
graph.  For a matroid $\mathcal M$ of rank $r$ the top $h$-number is
$T_{\mathcal M}(0,1)$, and $T_{\mathcal M^*}(x,y)=T_{\mathcal M}(y,x)$, where $T$ is
the Tutte polynomial.
\end{enumerate}
\item\label{I:cuts} \emph{Cuts of a graph.}  Let $\Gamma$ be a connected graph,
possibly with parallel edges, with one vertex $\infty$ marked.  For a set $I$
of vertices, the \emph{cut} $\delta(I)$ is the set of edges with exactly one
end in $I$.  A \emph{bond} is a minimal nonempty cut.
\begin{enumerate}[label=(\alph*)]
\item Give the edge between $D$ and $E$ the vector $x_D-x_E$, and the edge
between $D$ and $\infty$ the vector $x_D$.  This is the usual representation of
the graphic matroid, with $x_\infty=0$.  Its rank is the number of vertices
other than $\infty$, its cocircuits are the bonds, and every bond is
$\delta(I)$ for some $I$ not containing $\infty$.
\begin{details}
A set of edges spans the whole space exactly when it
connects all the vertices.  So the complement of a set $C$ of edges spans a
proper subspace exactly when deleting $C$ disconnects $\Gamma$, and the
cocircuits are the minimal such $C$, which are the bonds.  For $I$ take the
side of the disconnection without $\infty$.
\end{details}
\item Every cut is a disjoint union of bonds.
\begin{details}
If $S=\delta(I)$ is
nonempty, it contains a bond $B$, a minimal nonempty cut inside $S$.  Cuts are
closed under symmetric difference, since
$\delta(I)\,\triangle\,\delta(I')=\delta(I\,\triangle\,I')$, so $S\setminus B$
is a cut with fewer edges, and we induct.
\end{details}
\item The vectors of~(a) are unimodular \cite[Lemma~8.1]{hauselToricHyperKahler2002}.
\item The number of acyclic orientations of $\Gamma$ with $\infty$ the only
source is $T_\Gamma(1,0)$ \cite{greeneInterpretationWhitney1983}, where $T_\Gamma$ is the
Tutte polynomial.
\end{enumerate}
\item\label{I:alg} \emph{Finite groups and graded algebras.}
\begin{enumerate}[label=(\alph*)]
\item Let a finite group $H$ act linearly on a complex vector space $U$.  The
average $\Av_H=|H|^{-1}\sum_{h\in H}h$ is an $H$-equivariant projection of $U$
onto $U^H$.  So for an $H$-stable subspace $U'\subseteq U$ we have
$\Av_H(U')=U'^H$ and $(U/U')^H=U^H/U'^H$.
\item The symmetric polynomials in $n$ variables are generated by the power
sums of degrees $1,\dots,n$ (Newton's identities).
\item Let $C=\bigoplus_{i\ge0}C_i$ be a finite-dimensional graded commutative
algebra with $C_0=\C$, and $C_+=\bigoplus_{i>0}C_i$.  The \emph{socle} of $C$
is the annihilator of $C_+$, and $C$ is \emph{level} if its socle is its top
nonzero degree.  Every nonzero graded ideal $K$ of $C$ meets the socle.
\begin{details}
If $x\in K$ is homogeneous and nonzero of the largest degree
occurring in $K$, then $C_+x\subseteq K$ lies in higher degrees, so
$C_+x=0$.
\end{details}
\end{enumerate}
\item\label{I:sylv} \emph{Sylvester's equation.}  Let $A\in\mathrm{End}(U)$
and $B\in\mathrm{End}(U')$ have disjoint spectra.  Then $X\mapsto AX-XB$ is an
invertible map of $\Hom(U',U)$, and the solution of $AX-XB=C$ is a linear
combination of maps $A^rCB^s$.  Similarly, if $U=\bigoplus_DU_D$ and
$A=\bigoplus_DA_D$ with pairwise disjoint spectra, then the projection of
$U$ onto each $U_D$ is a polynomial in $A$.
\begin{details}
Let $\chi_B(t)=\sum_nc_nt^n$ be the characteristic polynomial of $B$.  The
eigenvalues of $\chi_B(A)$ are the $\chi_B(\alpha)$ for the eigenvalues
$\alpha$ of $A$, which are nonzero, so $\chi_B(A)$ is invertible, and by
Cayley--Hamilton its inverse is a polynomial in $A$.  For
$X\in\Hom(U',U)$ we have
$A^nX-XB^n=\sum_{i+j=n-1}A^i(AX-XB)B^j$, and $\chi_B(B)=0$, so
\[
 X=\chi_B(A)^{-1}\sum_nc_n\sum_{i+j=n-1}A^i\,(AX-XB)\,B^j .
\]
For the projections, by the Chinese remainder
theorem there is a polynomial $f_D$ that is $1$ modulo $\chi_{A_D}$ and $0$
modulo $\chi_{A_E}$ for $E\neq D$, and $f_D(A)$ is that projection.
\end{details}
\item\label{I:core} \emph{Cores} \cite{hauselCohomologyLarge2013}.  Let $X$ be a quiver variety
at a generic parameter $\nu$, as in \S\ref{ss:setup}.  The map
$\pi\colon X\to\mu^{-1}(0)/\!/G$ to the affine quotient is projective.  The
\emph{core} of $X$ is the projective variety $\fL=\pi^{-1}(\pi(0))$.  Let
$\C^\times$ scale $T^*N$ with weight one.  This preserves $\mu^{-1}(0)$ and
$\nu$-stability, so it acts on $X$, and for every $x\in X$ the limit
$\lim_{t\to0}t\cdot x$ exists in $X$ and lies in $\fL$.
\begin{details}
$\C^\times$ acts on the coordinate ring
$\C[\mu^{-1}(0)]^G$ of the affine quotient with positive weights apart from
the constants.  So $\pi(t\cdot x)\to\pi(0)$ as $t\to0$, and $\pi$ is
proper.
\end{details}
\begin{enumerate}[label=(\alph*)]
\item $X$ retracts onto $\fL$, so restriction $H^*(X)\to H^*(\fL)$ is an
isomorphism \cite{hauselCohomologyLarge2013}.
\item Every $G$-invariant function on $\mu^{-1}(0)$ that vanishes at $0$
vanishes on (the preimage of) $\fL$.  For example, for every loop $\ell$ of the
quiver, the traces of the powers of $\ell$ vanish on $\fL$, so $\ell$ is
nilpotent there.
\item If $p\in\mu^{-1}(0)$ is stable and $\lambda(t)\cdot p\to0$ as $t\to0$ for
a cocharacter $\lambda$ of $G$, then $[p]\in\fL$.
\begin{details}
Every $G$-invariant function $f$ satisfies $f(p)=f(\lambda(t)\cdot p)\to f(0)$, so
$\pi([p])=\pi(0)$.
\end{details}
\item If $\dim\fL=e$, then $\dim H^{2e}(\fL)$ is the number of irreducible
components of $\fL$ of dimension $e$ \cite[\S2.6]{CG97}.
\begin{details}
$\fL$ is compact, so
$H^{2e}(\fL)$ is dual to $H_{2e}(\fL)$, which has a basis formed by the
fundamental classes of these components \cite[\S2.6]{CG97}.
\end{details}
\end{enumerate}
\item\label{I:sheaves} \emph{Weights and the decomposition theorem}, used
only in \S\ref{sec:higgs}.  Complexes of sheaves on complex varieties are
taken in the derived category of mixed Hodge modules
\cite{saitoMixedHodge1990}, or, after spreading out, of mixed $\ell$-adic
complexes \cite{BBD}.  Perverse sheaves are normalized so that a local
system $\cL$ on a smooth variety $Y$ gives the perverse sheaf $\cL[\dim Y]$,
and $\IC(\overline Y,\cL)$ is its intermediate extension to $\overline Y$.
\begin{enumerate}[label=(\alph*)]
\item A complex $K$ has weights $\le w$ if and only if each stalk
$H^i(K_x)$ has weights $\le w+i$ (the definition \cite[\S5.1.8]{BBD}).  For
$f$ a morphism, $Rf_*$ preserves weights $\ge w$ \cite[\S5.1.14]{BBD}.  On a
smooth variety, $\Q$ is pure of weight $0$.  If $X$ is smooth and
$H^*(X)\to H^*(\fL)$ is an isomorphism for some proper $\fL\subseteq X$, then
$H^i(X)$ is pure of weight $i$: smoothness gives weights $\ge i$, and
properness of $\fL$ gives weights $\le i$.
\item A pure complex is, after extending scalars to an algebraic closure
in the $\ell$-adic case, a direct sum of shifted simple perverse sheaves
$\IC(\overline Y,\cL)$, with $Y$ smooth irreducible and $\cL$ an irreducible
local system on $Y$ \cite[\S\S5.3.8, 5.4.5]{BBD}, \cite{saitoMixedHodge1990}.
\item Let $y$ be a point of $\overline Y$ with $\dim Y>0$, such that
$\IC(\overline Y,\cL)$ is constructible for a stratification having $\{y\}$
as a stratum, and let $i_y$ be its inclusion.  Then
$H^j(i_y^*\IC(\overline Y,\cL))=0$ for $j\ge0$ and
$H^j(i_y^!\IC(\overline Y,\cL))=0$ for $j\le0$.  So the natural map
$i_y^!\to i_y^*$ is zero on it.
\item Let $\C^\times$ act on an affine space $V$ with positive weights, and
let $K$ be constructible for a stratification by $\C^\times$-stable strata.
Then restriction $R\Gamma(V,K)\to K_0$ is an isomorphism
\cite[Prop.~3.7.5]{kashiwaraSheavesManifolds1994}.
\begin{details}
The homeomorphism
$z\mapsto|z|^{1/a-1}z$ on each coordinate of weight $a$ turns the action of
$\mathbb R_{>0}$ into scaling, and $K$ is then conic.
\end{details}
\item Let a finite group $W$ act linearly on $\A^l$, with quotient
$q\colon\A^l\to\A^l/W$, and let $\cL(V)$ be the local system given by a
$W$-module $V$ on the image of the free locus.  Then
$\IC(\A^l/W,\cL(V))=(q_*\Q[l]\otimes V)^W$.  Its stalk at $q(0)$ is $V^W$ in degree $-l$, and the restriction
from this stalk to the stalk at a free point is the inclusion $V^W\subseteq
V$.
\end{enumerate}
\end{enumerate}

\section{The comparison map}\label{sec:map}
We regard $R_\hbar$ as the subalgebra of $A_{F,\hbar}$ of classes supported
over the base point.  For generic $\nu$ the map $R_\hbar\to B^\nu(A_{F,\hbar})$
is onto.  At $\hbar=0$ this is elementary, and we can say more:

\begin{lemma}\label{lem:generic}
Let $\nu$ be generic, let $A_{[k]}\subseteq A_F$ be the part of
$\nu$-degree $k$, and let $I_B\subseteq A_F$ be the ideal generated by the
elements of nonzero charge.  Put
\[
 J_\nu=\sum_{k>0}A_{[-k]}A_{[k]}.
\]
Then $A_{[0]}=R+J_\nu$ and $I_B\cap A_{[0]}=J_\nu$.  So $R\to
B^\nu_F$ is onto, and $B^\nu_F=B_F$.
\end{lemma}

\begin{proof}
The charge of a coweight $\lambda$ of $G$ is the vector of the sums of its
entries at each vertex.  So a generator $M_{\lambda,p}$ of~\ref{I:gen} with
$\lambda\neq0$ has a charge $c$ with $0<c\le\bv$ or $0<-c\le\bv$, and
$\deg_\nu(c)\neq0$ as $\nu$ is generic.  The generators with $\lambda=0$ are
the elements of $R$.  Now $A_{[0]}$ is spanned by the words in the generators
whose $\nu$-degrees add up to zero.  If such a word has a letter of nonzero
charge, then some letter has degree $-k<0$, and moving it to the front shows
that the word lies in $A_{[-k]}A_{[k]}$, as $A_F$ is commutative.  Otherwise
the word lies in $R$.  So $A_{[0]}=R+J_\nu$.  The same argument shows that
every charge component $A^{(c)}$ with $c\neq0$ and $\deg_\nu(c)=0$ lies in
$J_\nu$, since each word of charge $c$ has a letter of nonzero charge.

Decomposing $I_B$ by charge, this gives $I_B\cap A_{[0]}=J_\nu$.  As
$A_F=A^{(0)}+I_B$ and $A^{(0)}\subseteq A_{[0]}$, the map
$B^\nu_F=A_{[0]}/J_\nu\to A_F/I_B=B_F$ is an isomorphism.
\begin{details}
Clearly $J_\nu\subseteq I_B$.  Conversely, $I_B$ is graded by charge.  Its
components of nonzero charge $c$ with $\deg_\nu(c)=0$ lie in $A^{(c)}\subseteq
J_\nu$.  Its charge-zero part is $\sum_{c\neq0}A^{(-c)}A^{(c)}$, and each
$A^{(-c)}A^{(c)}$ lies in $J_\nu$: when $\deg_\nu(c)\neq0$ by the definition
of $J_\nu$, and otherwise because $A^{(c)}\subseteq J_\nu$, which is an ideal of
$A_{[0]}$.
\end{details}
\end{proof}

For $\hbar\neq0$ the algebra $A_{F,\hbar}$ is not commutative, and we reorder:

\begin{lemma}\label{lem:qsurj}
Let $\nu$ be generic, and let $A_{[k]}\subseteq A_{F,\hbar}$ be the part of
$\nu$-degree $k$.  Then
\[
 A_{[0]}=R_\hbar+J,\qquad J=\sum_{k>0}A_{[-k]}A_{[k]},
\]
so $R_\hbar\to B^\nu(A_{F,\hbar})$ is onto.
\end{lemma}

\begin{proof}
Fix a dominant coweight $\Lambda$.  The filtered piece $A_{\le\Lambda}$
of~\ref{I:gen} is a finitely generated graded $R_\hbar$-module, so the degrees
of $A_{\le\Lambda}\cap A_{[0]}$ are bounded below, and we show that every
homogeneous element of it lies in $R_\hbar+J$ by induction on degree.
By~\ref{I:gen}, such an element is a combination of words
$y=m_1\cdots m_r$ in the generators, of the same degree and of $\nu$-degree
zero, whose coweights sum to at most $\Lambda$.  If no letter has nonzero
charge, then $y\in R_\hbar$.  Otherwise, as in the proof of
Lemma~\ref{lem:generic}, some letter $m_j$ has $\nu$-degree $-k<0$, and
\[
 y=m_j\cdot m_1\cdots\widehat{m_j}\cdots m_r
 +\sum_{i<j}m_1\cdots m_{i-1}\,[m_i,m_j]\,m_{i+1}\cdots\widehat{m_j}\cdots m_r .
\]
The first term lies in $A_{[-k]}A_{[k]}$.  Since $A_F$ is commutative,
$[m_i,m_j]=\hbar c$ for some $c$, and $c$ lies in the same filtered piece as
$[m_i,m_j]$, because the subquotients of the filtration are free over
$\C[\hbar]$ \cite[\S6(i)]{BFN}.  So each other term is $\hbar$ times an
element of $A_{\le\Lambda}\cap A_{[0]}$ of lower degree, which lies in
$R_\hbar+J$ by induction.
\end{proof}

% BEN (original note, kept as a comment): I hope to convince Spencer and Tommaso
% to add this to their paper, otherwise, we can include the argument here.
The comparison map is the following.

\begin{theorem}[Botta--Tamagni, announced]\label{thm:map}
For every generic $\nu$, the surjective maps $R_{\hbar}\to B^\nu(A_{F,\hbar})$
and $R_{\hbar}\to H^*_{F\times\C^\times_\hbar}(X_\nu)$ induce a surjective
homomorphism of graded $R_\hbar$-algebras
\[
 \Phi\colon B^\nu(A_{F,\hbar})\to H^*_{F\times\C^\times_\hbar}(X_\nu).
\]
\end{theorem}

This is announced in \cite[Thm.~E, \S0.1.6.4]{bottaQuasimapCritical2026}, as a
consequence of the Coulomb branch action on the critical cohomology of based
quasimaps constructed there.  We use it without giving a proof.

Setting $\hbar=0$, and also the flavor parameters to zero, gives
surjections
\[
 \Phi\colon B^\nu_F\to H^*_F(X_\nu),\qquad\Phi\colon B^\nu_0\to H^*(X_\nu),
\]
with $\Phi[f]=\kappa(f)$ for $f\in R$.  Indeed, the B-algebras specialize
by~\ref{I:base}, and so does the equivariant cohomology, as it is free over
$H^*_{F\times\C^\times_\hbar}(\mathrm{pt})$ by~\ref{I:kirwan}.

\section{Quivers with a loop at every vertex}\label{sec:loops}

In this section $F$ is trivial, $\hbar=0$, every vertex $k$ with $v_k>0$
carries a loop, and $\nu=\nu_{\det}$.  Write $A=A_0$, $B=B_0$ and
$X=X_{\nu_{\det}}$, and let $I_B\subseteq A$ be the ideal generated by the
elements of nonzero charge, so that $B=A/I_B$.

As $\nu_{\det}$ is generic, Lemma~\ref{lem:generic} gives
$B^{\nu_{\det}}_0=B$, and $R\to B$ is onto.  By \S\ref{sec:map} we have a
surjection $\Phi\colon B\to H^*(X)$ with $\Phi[f]=\kappa(f)$ for $f\in R$.  We
grade by half the cohomological degree.

\begin{theorem}\label{thm:loops}
For every quiver with a loop at every vertex of positive dimension,
$\Phi\colon B\to H^*(X)$ is an isomorphism, and the Poincar\'e polynomial
$P_X(q)=\sum_i\dim H^{2i}(X)\,q^i$ of $X$ is
\[
 P_X(q)=\sum_mq^{\delta_m}h_m(q),
\]
the sum over multipartitions $m$, with $\delta_m$ and $h_m$ as in
\S\ref{ss:notation}.
\end{theorem}

\pre{Prop.~6.4} derives this formula for $P_X$ from Hausel's
formula \cite{hauselBettiNumbers2006}.  Here it comes out of the proof;
Remark~\ref{rem:hausel} explains how the two are related, and
Example~\ref{ex:A2loops} works out a small case.  For the Jordan
quiver with $\bw=1$, so that $X=\Hilb^n(\C^2)$, $Q'$ has no arrows, each
$\Pol_m/J_m$ is $\C$, and $\delta_m=n-l(m)$; the formula is
$P_X(q)=\sum_{m\vdash n}q^{n-l(m)}$.  For $\bw=r$, $X$ is the Gieseker variety
$M(r,n)$; now $\Gamma_m$ has $r$ edges from each part to $\infty$ and none
between parts, so $\Pol_m/J_m=\bigotimes_D\C[x_D]/(x_D^r)$ is the cohomology of
$(\mathbb P^{r-1})^{l(m)}$, and $\delta_m=r(n-l(m))$.  If $m$ has $a_s$ parts of
size $s$, then $\cW_m=\prod_s\mathfrak S_{a_s}$, the invariants are the
cohomology of $\prod_s\Sym^{a_s}\mathbb P^{r-1}$, and
\[
 P_{M(r,n)}(q)=\sum_{m\vdash n}q^{r(n-l(m))}\prod_s\binom{a_s+r-1}{a_s}_{\!q},
\]
with $q$-binomial coefficients.  For $r=2$ and $n=2$ this is
$(1+q+q^2)+q^2(1+q)=1+q+2q^2+q^3$.

\emph{Plan.}  We follow the outline in the introduction.
\begin{enumerate}[leftmargin=2em,label=(\alph*)]
\item \emph{The filtration} (\S\ref{ss:filt}).  Proposition~\ref{prop:filt}
filters $B$ by subspaces $F_L$, indexed by the number $L$ of parts of a
multipartition, with surjections
$\bigoplus_{l(m)=L}(\Pol_m/J_m)^{\cW_m}\langle\delta_m\rangle\twoheadrightarrow
F_L/F_{L-1}$, where $\Pol_m/J_m$ is the cohomology of a hypertoric variety
(Proposition~\ref{prop:toric}).  So $\dim B\le H(Q,\bv,\bw)$.
\item \emph{The induced filtration} (\S\ref{ss:induced}).  The images
$F'_L=\Phi(F_L)$ filter $H^*(X)$, and it suffices to show that the map from
$\bigoplus_{l(m)=L}(\Pol_m/J_m)^{\cW_m}\langle\delta_m\rangle$ to
$F'_L/F'_{L-1}$ is injective (Lemma~\ref{lem:reduction}).
\item \emph{Levi subgroups} (\S\ref{ss:levi}).  Restriction to the Levi
subgroup $M_m$, on the Coulomb side, and pullback along an \'etale map
$X(M_m)^\circ\to X$, on the quiver variety side, turn the summand of $m$ into
the piece $\Pol_m/J_m$ of the multipartition with one part per block for the
quiver $Q_{M_m}$, with no invariants taken (Proposition~\ref{prop:detect}).
So it suffices to treat that piece for every quiver with a loop at every
vertex.
\item \emph{The socle and the top degree} (\S\ref{ss:top}).  That piece is a
level module over $R$, so it suffices to check injectivity on its socle, its
top degree (Lemma~\ref{lem:socle}).  There we need enough classes in the top
degree of $H^*(X)$, and we construct by hand enough components of the core of
$X$ (Lemma~\ref{lem:comp}, Corollary~\ref{cor:topclass}).
\end{enumerate}
\S\ref{ss:proof} puts these together.

\subsection{Notation}\label{ss:notation}
Let $Q_0^+=\{k:v_k>0\}$ and $k_0=|Q_0^+|$.  Choose a loop $\ell_k$ at each
$k\in Q_0^+$, and let $Q'$ be $Q$ with these loops removed.  Put
$N'=N_{Q'}\oplus\Hom(W,V)$, so $N=\mathfrak g\oplus N'$.  Let $c_{kl}$ be the
number of arrows of $Q'$ between $k$ and $l$, in either direction, so that a
loop of $Q'$ at $k$ counts twice in $c_{kk}$.  Put
$z=\frac12\sum_kc_{kk}v_k$, the number of zero weights of the diagonal
torus on $N'$.  A point of $T^*N$ has two components in $\mathrm{End}(V_k)$
for the chosen loop $\ell_k$: the map along the loop, which we also call
$\ell_k$, and the map along the reversed loop, coming from the doubling of
the quiver, which we call $\ell_k^*$.  So $\ell_k\in\mathfrak{gl}(V_k)$ is the
coordinate on $N$, and $\ell_k^*$ is the dual coordinate on $N^*$.

\emph{Multipartitions and Levi subgroups.}  A \emph{multipartition} $m$ is
a tuple of partitions
\[
 m^{(k)}=(m^{(k)}_1\ge\cdots\ge m^{(k)}_{r_k})
\]
of $v_k$, $k\in Q_0^+$, where $r_k$ is the length of $m^{(k)}$.  Let
$P_m=\{(k,i):k\in Q_0^+,\ 1\le i\le r_k\}$, the set indexing the parts of $m$.  For $D=(k,i)\in P_m$ put $k(D)=k$ and
$s_D=m^{(k)}_i$, and put $l(m)=|P_m|$; we call the elements of $P_m$ the
\emph{parts} of $m$.
A decomposition
\[
 V_k=\bigoplus_{D\in P_m,\ k(D)=k}V^{(D)},\qquad \dim V^{(D)}=s_D,
\]
gives a Levi subgroup $M_m=\prod_{D\in P_m}\GL(V^{(D)})$.
Every Levi subgroup of $G$ is conjugate to $M_m$ for exactly one $m$.  

\emph{Toric data.}  Let
$\Zm\cong(\C^\times)^{l(m)}$ be the center of $M_m$.  The group
$\cW_m=N_G(M_m)/M_m$ permutes the parts of equal size at each vertex.  Let
$\xi_m$ be a regular nilpotent element of $\mathrm{Lie}(M_m)$, so that as an
element of $\mathfrak g$ it has Jordan type $m$.  The multipartition $m_G$,
with one part at each vertex, has $M_{m_G}=G$.

Let $N'^{\xi_m}\subseteq N'$ be the subspace killed by $\xi_m$.  The torus
$\Zm$ commutes with $\xi_m$, so it acts on $N'^{\xi_m}$.  Write
$x_D\in H^2_{\Zm}(\mathrm{pt})$ for the character by which $\Zm$ acts on
$V^{(D)}$, so that $\Pol_m=H^*_{\Zm}(\mathrm{pt})=\C[x_D:D\in P_m]$.  The space
$N'^{\xi_m}$ is the sum of the following pieces, on each of which $\Zm$ acts
by a single weight.
\begin{itemize}[leftmargin=2em]
\item For each arrow from $k$ to $l$ of $Q'$, and parts $D$ at $k$ and $E$ at
$l$, the space $\Hom(V^{(D)},V^{(E)})^{\xi_m}$ of maps commuting with
$\xi_m$, of dimension $\min(s_D,s_E)$, with weight $x_E-x_D$.  The weight is
zero exactly when $D=E$, for a loop of $Q'$.
\begin{details}
As a $\C[y]$-module, with $y$ acting by
$\xi_m$, each $V^{(D)}$ is $\C[y]/y^{s_D}$, so
\[
 \dim \Hom(V^{(D)},V^{(E)})^{\xi_m}=\dim \Hom_{\C[y]}(\C[y]/y^{s_D},\C[y]/y^{s_E})=\min(s_D,s_E).
\]
\end{details}
\item For each part $D$, the space $\Hom(W_{k(D)},V^{(D)}\cap\ker\xi_m)$, of
dimension $w_{k(D)}$, with weight $x_D$.
\begin{details}
$\Hom(W,V)^{\xi_m}=\Hom(W,\ker\xi_m)$, and $\ker\xi_m$ has one dimension in
each $V^{(D)}$.
\end{details}
\end{itemize}

\begin{definition}\label{def:toric-data}
Let $\Gamma_m$ be the graph whose vertices are the parts of $m$ and one more
vertex $\infty$.  It has $c_{k(D)k(E)}\min(s_D,s_E)$ edges between parts
$D\neq E$, and $w_{k(D)}$ edges between $D$ and $\infty$.  Label the edges
between $D$ and $E$ by $x_D-x_E$, and those between $D$ and $\infty$ by $x_D$.
Up to sign, these labels are the nonzero $\Zm$-weights of $N'^{\xi_m}$, one for
each weight vector.  For a nonempty set $I$ of parts, let $q_I\in \Pol_m$ be the
product of the labels of the edges with exactly one end in $I$, and let
$J_m\subseteq \Pol_m$ be the ideal generated by the $q_I$.  Put
\[
 h_m=\Hilb\big((\Pol_m/J_m)^{\cW_m}\big),\qquad
 \delta_m=\dim N'-\dim N'^{\xi_m},
\]
where $\Hilb$ is the Hilbert series in $q$, with each $x_D$ in degree $1$.
\end{definition}
For a graded space $U$, $U\langle\delta\rangle$ is $U$ with all degrees raised by
$\delta$.

\begin{proposition}\label{prop:toric}
We have a ring isomorphism $\Pol_m/J_m\cong H^*(\fM_m)$, where $\fM_m$ is the hypertoric variety
of $\Zm$ acting on $N'^{\xi_m}$ with its zero weights removed, at a generic
parameter.  It vanishes above degree $n_m-l(m)$, where $n_m$ is the number
of edges of $\Gamma_m$, and $\delta_m+n_m-l(m)=\dim N'-z-l(m)$.  If it is
nonzero, it is level: its socle is its top degree $n_m-l(m)$.  That degree
has dimension the number of acyclic orientations of $\Gamma_m$ in which
$\infty$ is the only source.
\end{proposition}

\begin{proof}
By~\ref{I:cuts}(a), with $x_\infty=0$, the weights of $\Zm$ on $N'^{\xi_m}$ are
the vectors of the graphic matroid of $\Gamma_m$.  Suppose first that $\Gamma_m$
is connected.  Then they span $H^2_{\Zm}(\mathrm{pt})$, and they are unimodular
by~\ref{I:cuts}(c).  By~\ref{I:hypertoric}(b), $H^*(\fM_m)$ is $\Pol_m$ modulo the
products of the weights over the cocircuits.  By~\ref{I:cuts}(a) the cocircuits
are the bonds, and each bond is $\delta(I)$ for a nonempty set $I$ of parts.  So
the ideal is contained in $J_m$.  Conversely, each $q_I$ is the product over the
cut $\delta(I)$, which by~\ref{I:cuts}(b) is a disjoint union of bonds, so
$q_I$ is a multiple of a bond product.  So the ideal is $J_m$.  (This is also
\cite[Thm.~8.3]{hauselToricHyperKahler2002}, which presents the cohomology of a
toric quiver variety by the products over all cuts.)  If $\Gamma_m$ is
disconnected, then a component $I$ not containing $\infty$ has $q_I=1$, so
$\Pol_m/J_m=0$.  The weights do not span, and a generic $\theta$ lies outside
their span, so $\fM_m$ is empty and the isomorphism still holds.

By~\ref{I:hypertoric}(a), $\fM_m$ is smooth, since the weights are
unimodular, and it has dimension $2(n_m-l(m))$ when it is nonempty.
By~\ref{I:hypertoric}(c), with $\dim\Zm=l(m)$ and $n_m$ nonzero weights, the
ring vanishes above degree $n_m-l(m)$, and it is level.

For the formula, the edges of $\Gamma_m$ correspond to the weight vectors of
$N'^{\xi_m}$ with nonzero weight, and there are $z$ zero weights, so
$n_m=\dim N'^{\xi_m}-z$ and $\delta_m+n_m=\dim N'-z$.
\begin{details}
By the description of $N'^{\xi_m}$
before Definition~\ref{def:toric-data}, the zero weights are those of
$\Hom(V^{(D)},V^{(D)})^{\xi_m}$ for the loops of $Q'$: each loop at $k$ and
each part $D$ at $k$ give $s_D$ of them, so there are $\sum_k\frac12c_{kk}v_k=z$.
\end{details}
By~\ref{I:hypertoric}(d) the top Betti number
is the top $h$-number of the cographic matroid of $\Gamma_m$, which is
$T_{\Gamma_m}(1,0)$.
By~\ref{I:cuts}(d) this counts the acyclic orientations of $\Gamma_m$ in which
$\infty$ is the only source.
\end{proof}

\begin{remark}[the cuts as Euler classes]\label{rem:kirwan}
The ideal $J_m$ has a direct geometric meaning, which we learned from
\cite[\S\S1--2]{nakajimaMultipartitionSpaces2026}, notes on \pre{\S6} written by Claude at the request of
H.~Nakajima.  Let
$E_m$ be the space of representations, with dimension one at every vertex, of
the quiver with vertices $P_m\sqcup\{\infty\}$, one arrow in each direction for
each edge of $\Gamma_m$ between two parts, and one arrow $\infty\to D$ for each
edge between $D$ and $\infty$.  The torus $\Zm$ acts on $E_m$, trivially at
$\infty$, with weight $x_E-x_D$ on an arrow $D\to E$ and $x_D$ on an arrow
$\infty\to D$.  Let $U_m\subseteq E_m$ be the open set where every part can be
reached from $\infty$ along nonzero arrows.  As in~\ref{I:cyc}, it is the
semistable locus for the determinant character, and $\Zm$ acts freely on it,
so $U_m/\Zm$ is a smooth toric variety.
\begin{enumerate}[leftmargin=2em]
\item Restriction $\Pol_m=H^*_{\Zm}(E_m)\to H^*_{\Zm}(U_m)=H^*(U_m/\Zm)$ is
onto, with kernel $J_m$.  This is Kirwan's method
\cite{kirwanCohomologyQuotients1984} for the torus $\Zm$ and the determinant
character, applied to the stratification of $E_m$ by the set $I$ of parts that
cannot be reached; the class of the stratum of $I$ is $\pm q_I$.
\begin{details}
The stratum of $I$ is open in the linear subspace $L_I$
where the arrows into $I$ from $\infty$ and from the parts outside $I$ vanish.
There is one such arrow for each edge of the cut $\delta(I)$, so the
equivariant class of $L_I$, which is the Euler class of the normal bundle of
the stratum, is $\pm q_I$.  The stratum retracts onto $U_{m'}$, where $m'$ is
the multipartition formed by the parts outside $I$, so its equivariant
cohomology is $H^*_{Z_{m'}}(U_{m'})\otimes\C[x_D:D\in I]$.  On this ring $q_I$
is not a zero divisor, since its part of top degree in the $x_D$, $D\in I$, is
a single monomial.  So the Thom--Gysin sequences for removing the strata in
order of increasing $|I|$ split; this is the criterion of Atiyah and Bott.  By
induction on the number of parts, every class on a stratum is restricted from
$\Pol_m$, so the image of its Gysin map is generated by $q_I$, and the kernel is
$J_m$.  These are Kirwan's strata: the optimal destabilizing
one-parameter subgroup of a point whose unreachable set is $I$ acts with
weight $-1$ on the parts in $I$ and trivially on the others.
\end{details}
\item Adding arrows $D\to\infty$ dual to the framing arrows turns $E_m$ into
the cotangent space of the representation space of an orientation of
$\Gamma_m$.  The new arrows do not help to reach parts from $\infty$, so the
quotient of its semistable locus by $\Zm$, which is the Lawrence toric variety
of \cite{hauselToricHyperKahler2002}, is a vector bundle over $U_m/\Zm$.  For the
determinant character it contains $\fM_m$ as the zero set of the moment map,
and $\fM_m$ is a deformation retract of it \cite[Remark after the proof of
Thm.~1.1]{hauselToricHyperKahler2002}.  So~(1) gives the ring isomorphism of
Proposition~\ref{prop:toric} again, with all cuts at once.  In this
description, $\fM_m$ is the quiver variety of an orientation of $\Gamma_m$,
with dimension one at each part and the edges at $\infty$ as framing, at
determinant stability.
\item Compare Proposition~\ref{prop:filt}(3) below.  There, the cut relation
for $I$ comes from splitting off the coweight supported on the first elements
of the parts in $I$, whose nonzero charge kills it in $B$, and the coefficient
of $r_{\lambda_m}$ in the resulting product is $\pm q_I$.  Here $q_I$ is the
class of the stratum where exactly the parts in $I$ are unreachable.
\item The group $\cW_m$ permutes the strata, and taking invariants in the split
Thom--Gysin sequences gives
\[
 \frac1{(q)_m}=\sum_I q^{\deg q_I}\,\frac{h_{m'}(q)}{(q)_I},
\]
the sum over the $\cW_m$-orbits of subsets $I\subseteq P_m$, where $m'$ is the
multipartition formed by the parts outside $I$.  Here $1/(q)_m=\prod_{k,s}\prod_{j=1}^{a_{k,s}}(1-q^j)^{-1}$ is the
Hilbert series of $\Pol_m^{\cW_m}$, where $a_{k,s}$ is the number of parts of
$m$ of size $s$ at $k$, and $(q)_I$ is defined in the same way from the parts
in $I$.  The term of $I=\emptyset$ is $h_m$, so this determines $h_m$
recursively.  It is \pre{Prop.~6.2}; see \cite[Prop.~2.1]{nakajimaMultipartitionSpaces2026}.
\end{enumerate}
Nothing in this remark is used in the proof of Theorem~\ref{thm:loops}.
\end{remark}

\begin{example}\label{ex:A2loops}
Let $Q$ be the quiver $1\to2$ with a loop added at each vertex, so that $Q'$
is $1\to2$, $c_{12}=1$ and $c_{11}=c_{22}=0$, and take $\bv=(2,1)$ and
$\bw=(1,0)$.  Then $\dim N'=v_1v_2+w_1v_1=4$, $z=0$ and $k_0=2$.  There are
two multipartitions.
\begin{itemize}[leftmargin=2em]
\item $m_G=((2),(1))$, with a part $D$ of size $2$ at $1$ and a part $E$ of
size $1$ at $2$.  The graph $\Gamma_{m_G}$ is the path $\infty-D-E$, as
$w_1=1$, $w_2=0$ and $\min(2,1)=1$.  Up to sign, $q_{\{D\}}=x_D(x_D-x_E)$,
$q_{\{E\}}=x_E-x_D$ and $q_{\{D,E\}}=x_D$, so $J_{m_G}=(x_D,x_E)$ and
$h_{m_G}=1$.  Also $\dim N'^{\xi_{m_G}}=\min(2,1)+w_1=2$, so $\delta_{m_G}=2$.
\item $m=((1,1),(1))$, with parts $D_1,D_2$ at $1$ and $E$ at $2$, all of size
one.  Then $\xi_m=0$, so $\delta_m=0$, and $\Gamma_m$ is the $4$-cycle
$\infty-D_1-E-D_2-\infty$.  With $a=x_{D_1}$, $b=x_{D_2}$ and $c=x_E$, the
$q_I$ are, up to sign, $a(a-c)$, $b(b-c)$, $(c-a)(c-b)$, $a(c-b)$, $b(c-a)$,
$ab$ and $ab(a-c)(b-c)$.  The first six generate $(a,b,c)^2$, so
$J_m=(a,b,c)^2$, and $\Pol_m/J_m$ has basis $1,a,b,c$.  Its top degree
$n_m-l(m)=1$ has dimension $3$, the number of acyclic orientations of the
$4$-cycle with $\infty$ the only source, and $\fM_m$ is the minimal resolution
of $\C^2/(\Z/4)$.  The group $\cW_m=\mathfrak S_2$ swaps $a$ and $b$, so
$h_m=1+2q$.
\end{itemize}
By Theorem~\ref{thm:loops}, the $8$-dimensional quiver variety $X$ has
$P_X(q)=(1+2q)+q^2$, with top degree $d_{\max}=\dim N'-z-k_0=2$.  This example is taken from
\cite[Exs.~1.13, 3.8]{nakajimaMultipartitionSpaces2026}.
\end{example}

\subsection{The Coulomb branch in abelianized form}\label{ss:abel}
Let $T\subseteq G$ be the diagonal torus.  Index the standard bases of the
$V_k$ by $\Omega=\{(k,a):k\in Q_0,\ 1\le a\le v_k\}$, and write $\Omega_k$ for
the elements with first entry $k$, so $|\Omega_k|=v_k$.  So
$T=(\C^\times)^\Omega$, coweights are elements of $\Z^\Omega$, and the
charge of $\lambda$ is $(\sum_{d\in\Omega_k}\lambda_d)_k$.  Also
$\Pol=H^*_T(\mathrm{pt})=\C[x_d:d\in\Omega]$, and $\cW=\prod_k\mathfrak S_{v_k}$.  Let
$\cB'$ be the multiset of $T$-weights of $N'$.  In the commutative ring
$\mathrm{Frac}(\Pol)[\Z^\Omega]$, with basis $e^\lambda$ over
$\mathrm{Frac}(\Pol)$, put
\[
 r_\lambda=e^\lambda\prod_{\beta\in\cB'}\beta^{(-\langle\beta,\lambda\rangle)_+},
\]
where $y_+=\max(y,0)$ and $\langle\ ,\ \rangle$ pairs weights with coweights.
The group $\cW$ acts on this ring by permuting the $x_d$ and the coordinates
of $\lambda$ together.

\begin{proposition}[Teleman]\label{prop:teleman}
$A=\big(\bigoplus_{\lambda\in\Z^\Omega}\Pol\,r_\lambda\big)^{\cW}$, compatibly
with the abelianization embedding of \cite[Lemma~5.17]{BFN}.
\end{proposition}

This is \cite[App.~A, Thm.~6]{telemanRoleCoulomb2021}: if the weights of $N$ contain
the roots, the Coulomb branch is the $\cW$-invariant part of the abelian
Coulomb branch of $T$ acting on $N$ with the roots removed.  The chosen
loops supply the roots: the weights of $\ell_k$ are the $x_e-x_d$ for
$d,e\in\Omega_k$, so removing the roots leaves the weights of $N'$ and some
zero weights, which do not contribute to $r_\lambda$.

For $\lambda\in\Z^\Omega$, let $\cW_\lambda$ be its stabilizer in $\cW$.
Since $\cB'$ is $\cW$-stable, $w(r_\lambda)=r_{w\lambda}$ for $w\in\cW$.  For
$p\in \Pol^{\cW_\lambda}$, the \emph{dressed
monopole} $\mathsf S_\lambda(p)=\sum_{w\in\cW/\cW_\lambda}w(p\,r_\lambda)$ lies
in $A$ and has the charge of $\lambda$.  Products are computed with
\begin{equation}\label{eq:rr}
 r_\lambda r_\mu=\Big(\prod_{\beta\in\cB'}\beta^{(-\langle\beta,\lambda\rangle)_+
 +(-\langle\beta,\mu\rangle)_+-(-\langle\beta,\lambda+\mu\rangle)_+}\Big)\,r_{\lambda+\mu},
\end{equation}
whose prefactor is a polynomial.  Every element of $A$ is a sum of dressed
monopoles, one for each $\cW$-orbit of coweights.
\begin{details}
Each $\cW$-invariant element
$\sum_\lambda f_\lambda r_\lambda$, with $f_\lambda\in\Pol$, has
$f_{w\lambda}=w(f_\lambda)$, so $f_\lambda\in\Pol^{\cW_\lambda}$, and its terms
on the $\cW$-orbit of $\lambda$ add up to $\mathsf S_\lambda(f_\lambda)$.
\end{details}

\emph{Degrees.}  We grade by half the cohomological degree, so each $x_d$
has degree $1$.  In the grading of \cite{BFN}, halved, the class
$r_\lambda$ has degree $\frac12\sum_{\beta\in\cB'}|\langle\beta,\lambda\rangle|$.
If $\lambda$ has charge zero, this is
\[
 \deg r_\lambda=\sum_{\beta\in\cB'}(-\langle\beta,\lambda\rangle)_+,
\]
the degree of the prefactor of $r_\lambda$.
\begin{details}
The character $\sum_{\beta\in\cB'}\beta$ is $\cW$-invariant, hence a
combination of the characters $\sum_{d\in\Omega_k}x_d$.  So if $\lambda$ has
charge zero, then $\sum_{\beta\in\cB'}\langle\beta,\lambda\rangle=0$, and
$\frac12|y|=(-y)_++\frac12y$.
\end{details}

\subsection{The multipartition filtration}\label{ss:filt}
Fix a multipartition $m$.  Split each $\Omega_k$ into blocks, one for each
part $D$ at $k$, of size $s_D$, and choose a \emph{first element} $d_D$ in
each block.  Let $s\in\Z^\Omega_{\ge0}$ be $s_D$ at $d_D$ and $0$ elsewhere,
and put $\lambda_m=s-\one$, where $\one\in\Z^\Omega$ has all entries $1$.  So
$\lambda_m$ is $s_D-1$ at the first element of
the block $D$ and $-1$ at its other elements.  Listing the elements of a
block with its first element first, the entries of $\lambda_m$ on $D$ are
\[
 \lambda_m|_D=(\,\underbrace{s_D-1,\,-1,\,-1,\,\dots,\,-1}_{s_D\text{ entries}}\,),
\]
and $\lambda_m$ is the concatenation of these vectors over all blocks $D$.  For
example, for $m=(3,1)$ at a vertex with $v_k=4$ it is $(2,-1,-1,\,0)$.  It has
charge zero, and its
$\cW$-orbit does not depend on the choices.  The blocks also realize $M_m$,
with $T\subseteq M_m$.  Let $\epsilon\colon \Pol\to \Pol_m$ be restriction to
$\Zm\subseteq T$, which sets $x_d=x_D$ for $d$ in the block $D$.  For
$m=m_G$ we write $\lambda_G$, $\Pol_G$ and $J_G$.

Let $F_L\subseteq B$ be the span of the classes $[\mathsf S_{\lambda_m}(p)]$
with $l(m)\le L$ and $p\in \Pol^{\cW_{\lambda_m}}$.  It does not depend on the
choice of blocks and first elements: another choice replaces $\lambda_m$ by
some $w\lambda_m$, and $\mathsf S_{w\lambda}(wp)=\mathsf S_\lambda(p)$.

\begin{proposition}\label{prop:filt}
\begin{enumerate}[leftmargin=2em]
\item $F_{|\bv|}=B$ and $F_{k_0-1}=0$.
\item For $p\in\Pol^{\cW_{\lambda_m}}$, the class $[\mathsf S_{\lambda_m}(p)]\in B$
depends only on $\epsilon(p)$, and it has degree $\deg p+\delta_m$.
\item If $\epsilon(p)\in J_m$, then $[\mathsf S_{\lambda_m}(p)]\in
F_{l(m)-1}$.
\end{enumerate}
So $\epsilon(p)\mapsto[\mathsf S_{\lambda_m}(p)]$ induces a surjection
$\bigoplus_{l(m)=L}(\Pol_m/J_m)^{\cW_m}\langle\delta_m\rangle\twoheadrightarrow
F_L/F_{L-1}$.
\end{proposition}

\begin{proof}
(1) If every part of $m$ has size $1$, then every element of $\Omega$ is a
first element and $s=\one$, so $\lambda_m=0$, $\cW_{\lambda_m}=\cW$ and
$\mathsf S_0(p)=p$ for $p\in \Pol^\cW=R$.  So $F_{|\bv|}$ contains the image
of $R$, which is $B$, as shown at the start of this section.  Every
multipartition has at least one part at each vertex of $Q_0^+$, so
$l(m)\ge k_0$ and $F_{k_0-1}=0$.

(2) \emph{The stabilizer.}  Write $H=\cW_{\lambda_m}$.  It preserves the set
of first elements of each $\Omega_k$, on which it acts by the group $H_1$ of
permutations of the first elements of parts of equal size at the same vertex;
under $d_D\mapsto D$, $H_1$ is $\cW_m$.  On the non-first elements of each
$\Omega_k$ it acts by the full symmetric group, and $H$ is the product of $H_1$
and these symmetric groups.
\begin{details}
The entry of
$\lambda_m$ is $s_D-1\ge0$ at a first element $d_D$ and $-1$ at a non-first
element, and $H$ permutes exactly the elements of each $\Omega_k$ with equal
entries.
\end{details}

\emph{Generators of $\Pol^H$.}  Let $\Pol^{\rm f}=\C[x_{d_D}:D\in P_m]$ be the
polynomials in the first elements, and for $k\in Q_0^+$ and $j\ge1$ let
\[
 P_{k,j}=\sum_{d\in\Omega_k}x_d^j\in R,\qquad
 P^m_{k,j}=\sum_{D:\,k(D)=k}s_D\,x_{d_D}^j\in(\Pol^{\rm f})^{H_1}.
\]
By~\ref{I:alg}(b), $\Pol^H$ is generated over $(\Pol^{\rm f})^{H_1}$ by the
$P_{k,j}$.
\begin{details}
By~\ref{I:alg}(b), $\Pol^H$ is generated over $(\Pol^{\rm f})^{H_1}$ by the
power sums of the non-first variables of each $\Omega_k$, and the power sum of
the first variables of $\Omega_k$ lies in $(\Pol^{\rm f})^{H_1}$.
\end{details}

\emph{A relation from derivations.}  It suffices to show that multiplying
the dressing by $P_{k,j}$ or by $P^m_{k,j}$ gives the same class in $B$.  This
comes from a derivation of $A$ that acts by zero on $B$.  For each $k$ and
$j\ge1$, let
$\eta_{k,j}$ be the map
\[
 \eta_{k,j}\colon p\,e^\lambda\mapsto p\,e^\lambda\sum_{d\in\Omega_k}\lambda_dx_d^j
\]
on $\mathrm{Frac}(\Pol)[\Z^\Omega]$.  It is a derivation, since the multiplier
is additive in $\lambda$.  It commutes with $\cW$, since $\cW$ permutes each
$\Omega_k$, and it multiplies each $r_\lambda$ by a polynomial.  So it
preserves $A$ and its charge components, and as it is a derivation it
preserves the ideal $I_B$.  So it induces a map on $B=A/I_B$, and this map is
zero.  Indeed, $R\to B$ is onto, so every $a\in A$ is $r+i$ with $r\in R$ and
$i\in I_B$.  Now $\eta_{k,j}(r)=0$, since $r$ has $\lambda=0$, and
$\eta_{k,j}(i)\in I_B$.  So $[\eta_{k,j}(a)]=0$ in $B$ for every $a\in A$.

Take $a=\mathsf S_{\lambda_m}(p)=\sum_{w\in\cW/H}w(p\,r_{\lambda_m})$ with
$p\in\Pol^H$.  The multiplier of $\eta_{k,j}$ at $w\lambda_m$ is
$w(\sigma_{k,j})$, where $\sigma_{k,j}$ is its value at $\lambda_m$:
\[
 \sigma_{k,j}=\sum_{d\in\Omega_k}(\lambda_m)_dx_d^j
 =-P_{k,j}+P^m_{k,j}.
\]
It is $H$-invariant, so $\eta_{k,j}(a)=\mathsf S_{\lambda_m}(\sigma_{k,j}\,p)$.
Since $[\eta_{k,j}(a)]=0$, for all $p\in\Pol^H$ we have
\begin{equation}\label{eq:powersum}
 [\mathsf S_{\lambda_m}(P_{k,j}\,p)]=[\mathsf S_{\lambda_m}(P^m_{k,j}\,p)].
\end{equation}

\emph{Dependence on $\epsilon(p)$.}  Let $\epsilon^{\rm f}\colon\Pol^{\rm f}\to\Pol_m$ be the restriction of
$\epsilon$, the isomorphism with $x_{d_D}\mapsto x_D$.  The polynomials
$P_{k,j}$ and $P^m_{k,j}$ have the same image $\sum_Ds_Dx_D^j$ under
$\epsilon$, so applying~\eqref{eq:powersum} repeatedly shows that $[\mathsf
S_{\lambda_m}(p)]$ depends only on $\epsilon(p)$.  Also
$\epsilon(\Pol^H)=\Pol_m^{\cW_m}$.
\begin{details}
Put
$\varphi=(\epsilon^{\rm f})^{-1}\circ\epsilon$ on $\Pol^H$.  It is the algebra
homomorphism $\Pol^H\to(\Pol^{\rm f})^{H_1}$ that fixes $(\Pol^{\rm f})^{H_1}$
and sends $P_{k,j}$ to $P^m_{k,j}$, since both have image
$\sum_Ds_Dx_D^j$ under $\epsilon$.  Let $p=c\,P_{k_1,j_1}\cdots
P_{k_r,j_r}$ with $c\in(\Pol^{\rm f})^{H_1}$.  Applying~\eqref{eq:powersum}
$r$ times gives
\[
 [\mathsf S_{\lambda_m}(p)]=[\mathsf S_{\lambda_m}(c\,P^m_{k_1,j_1}\cdots
 P^m_{k_r,j_r})]=[\mathsf S_{\lambda_m}(\varphi(p))].
\]
Such $p$ span $\Pol^H$, so $[\mathsf S_{\lambda_m}(p)]=[\mathsf
S_{\lambda_m}(\varphi(p))]$ for all $p\in\Pol^H$, and this depends only on
$\epsilon(p)$.  Also $\varphi(\Pol^H)=(\Pol^{\rm f})^{H_1}$, so
$\epsilon(\Pol^H)=\Pol_m^{\cW_m}$.
\end{details}

\emph{Degree.}  The element $\mathsf S_{\lambda_m}(p)$ is a sum of $\cW$-translates of
$p\,r_{\lambda_m}$, and $\lambda_m$ has charge zero, so by \S\ref{ss:abel}
its degree is $\deg p+\sum_{\beta\in\cB'}(-\langle\beta,\lambda_m\rangle)_+$.
Counting the weights of $N'$ arrow by arrow shows that this sum is
$\dim N'-\dim N'^{\xi_m}=\delta_m$.
\begin{details}
Write $s_d$ for the entry of $s$ at $d\in\Omega$, so $s_{d_D}=s_D$ and $s_d=0$
at non-first elements.  Then $\langle\beta,\lambda_m\rangle
=\langle\beta,s\rangle-\langle\beta,\one\rangle$.
\begin{itemize}[leftmargin=2em]
\item An arrow from $k$ to $l$ of $Q'$ contributes the weights $x_e-x_d$, for
$d\in\Omega_k$ and $e\in\Omega_l$, with
$(-\langle\beta,\lambda_m\rangle)_+=(s_d-s_e)_+$.  This is $0$ if $d$ is not
a first element.  It is $s_D$ if $d=d_D$ and $e$ is not a first element, and
$\Omega_l$ has $v_l-r_l$ non-first elements.  It is
$(s_D-s_E)_+=s_D-\min(s_D,s_E)$ if $d=d_D$ and $e=d_E$.  Since
$\sum_{k(D)=k}s_D=v_k$, the total is
\[
 v_k(v_l-r_l)+r_lv_k-\sum_{k(D)=k,\,k(E)=l}\min(s_D,s_E)
 =v_kv_l-\sum_{k(D)=k,\,k(E)=l}\min(s_D,s_E).
\]
\item The framing at $k$ contributes the weights $x_d$, $d\in\Omega_k$, each
$w_k$ times, with $(-\langle\beta,\lambda_m\rangle)_+=(1-s_d)_+$.  This is $1$
at the $v_k-r_k$ non-first elements and $0$ at the first elements, so the
total is $w_k(v_k-r_k)$.
\end{itemize}
By the computation before Definition~\ref{def:toric-data},
\[
 \dim\Hom(V_k,V_l)^{\xi_m}=\sum_{k(D)=k,\,k(E)=l}\min(s_D,s_E),\qquad
 \dim\Hom(W_k,V_k)^{\xi_m}=w_kr_k .
\]
So the sum is $\dim N'-\dim N'^{\xi_m}=\delta_m$.
\end{details}

(3) For each nonempty set $I$ of parts, we split $\lambda_m$ into two
coweights, one of nonzero charge, and expand their product, which vanishes in
$B$.  Its leading term involves $q_I$, and the other terms lie in
$F_{l(m)-1}$.

\emph{The cut relation.}  Fix a nonempty set $I\subseteq P_m$ of parts.  Let $b\in\Z^\Omega_{\ge0}$
agree with $s$ at the first elements of the parts in $I$ and vanish
elsewhere, and put $a=s-b$.  So $a$ records the parts outside $I$, $b$ records
the parts in $I$, and $\lambda_m=(a-\one)+b$.  The charge of $b$ is nonzero,
since its entries are $\ge0$ and $s_D>0$ for $D\in I$.  So for
$p_a\in\Pol^{\cW_a}$ and $p_b\in\Pol^{\cW_b}$, the element
$\mathsf S_b(p_b)$ lies in $I_B$, and in $B$ we have
\begin{equation}\label{eq:cutrel}
 [\mathsf S_{a-\one}(p_a)\,\mathsf S_b(p_b)]=0.
\end{equation}
(Note $\cW_{a-\one}=\cW_a$.)  We expand this product.  Since
$w(r_\lambda)=r_{w\lambda}$,
\[
 \mathsf S_{a-\one}(p_a)\,\mathsf S_b(p_b)
 =\sum_{w\in\cW/\cW_a}\ \sum_{w'\in\cW/\cW_b}w(p_a)\,w'(p_b)\,
 r_{wa-\one}\,r_{w'b},
\]
and by~\eqref{eq:rr}, the term of $(w,w')$ is a polynomial times $r_{t-\one}$,
where $t=wa+w'b$.  By \S\ref{ss:abel}, grouping the terms by the
$\cW$-orbit of $t-\one$ writes the product as a sum of dressed monopoles, one
for each orbit.

\emph{Which orbits occur.}  Since $t\ge0$ and $\sum_{d\in\Omega_k}t_d=v_k$, the
nonzero entries of $t$ on each $\Omega_k$, in decreasing order, form a
partition of $v_k$.  Let $m'$ be the multipartition they define; then $t-\one$
is conjugate to $\lambda_{m'}$.  The supports of $a$ and $b$ are disjoint, with
$l(m)$ elements in all, so $l(m')$ is $l(m)$ minus the number of common
elements of the supports of $wa$ and $w'b$.  The orbits coming from terms
with overlapping supports therefore give dressed monopoles in
$F_{l(m)-1}$.  If the supports are disjoint, then the nonzero entries of $t$
at $k$ are those of $a$ and $b$ together, which are the parts of $m$ at $k$.
So $m'=m$, and all these terms lie on the orbit of $\lambda_m$.  That orbit
contributes $\mathsf S_{\lambda_m}(f)$, where $f$ is the coefficient of
$r_{\lambda_m}$, which comes from the pairs $(wa,w'b)$ with $wa+w'b=s$.

\emph{The coefficient of $r_{\lambda_m}$.}  These pairs are the $(ha,hb)$
for $h\in H$, and they are indexed by $H/H_I$, where $H_I=\cW_a\cap\cW_b$.
Write $r_{a-\one}r_b=g\,r_{\lambda_m}$, with $g$ the prefactor
of~\eqref{eq:rr}.  Since $p_a$ is $\cW_a$-invariant and $p_b$ is
$\cW_b$-invariant, the pair $(ha,hb)$ contributes $h(p_ap_b\,g)\,r_{\lambda_m}$,
so $f=\sum_{h\in H/H_I}h(p_ap_b\,g)$.  Comparing the exponents in $g$ with
Definition~\ref{def:toric-data}, $g=\pm\tilde q_I$, where $\tilde
q_I=(\epsilon^{\rm f})^{-1}(q_I)$ is $q_I$ written in the first elements.
\begin{details}
Given a pair $(wa,w'b)$ with $wa+w'b=s$,
$w'b$ agrees with $s$ on its support, and the first elements in that support
at each vertex carry the sizes of the parts of $I$ there.  So some $h\in H_1$
maps the first elements of the parts in $I$ onto them, preserving sizes, and
then $hb=w'b$ and $ha=h(s-b)=s-w'b=wa$.  The stabilizer of $(a,b)$ in $H$ is
$H_I$.

For a weight $x_e-x_d$ of an arrow from $k$ to $l$ of $Q'$,
its exponent in $g$ is
$(a_d-a_e)_++(b_d-b_e)_+-(s_d-s_e)_+$.  This is $0$ unless $d$ and $e$ are
both first elements, $d=d_D$ and $e=d_E$.  It is also $0$ if $D$ and $E$ are
both in $I$ or both outside $I$, and it is $\min(s_D,s_E)$ if exactly one of
them is in $I$.  For a framing weight $x_d$, the exponent is
$(1-a_d)_+-(1-s_d)_+$, which is $1$ if $d$ is the first element of a part in
$I$ and $0$ otherwise.
\end{details}
By~\eqref{eq:cutrel}, we have shown
\begin{equation}\label{eq:cutclass}
 \Big[\mathsf S_{\lambda_m}\Big(\textstyle\sum_{h\in H/H_I}h(\tilde
 q_I\,p_ap_b)\Big)\Big]\in F_{l(m)-1}.
\end{equation}

\emph{Choice of dressings.}  Averaging~\eqref{eq:cutclass} over suitable
dressings $p_a$ and $p_b$ gives $[\mathsf S_{\lambda_m}(\Av_H(\tilde
q_Iy))]\in F_{l(m)-1}$ for every $y\in\Pol^{\rm f}$, and the images
$\epsilon(\Av_H(\tilde q_Iy))=\Av_{\cW_m}(q_I\,\epsilon(y))$ span
$\Av_{\cW_m}(J_m)=J_m^{\cW_m}$.
\begin{details}
The group $\cW_a$ permutes the first elements
of the parts outside $I$ of equal size at the same vertex, and it permutes
all the other elements of each $\Omega_k$ arbitrarily, since $a$ vanishes
there.  So every polynomial $p_a$ in the variables $x_{d_D}$, $D\notin I$,
invariant under permuting the parts outside $I$ of equal size at the same
vertex, lies in $\Pol^{\cW_a}$.  Likewise for $p_b$, with $D\in I$.  On
$\Pol^{\rm f}$, the group $H_I$ acts by the permutations of the first elements
that preserve vertex and size and preserve $I$, and the products $p_ap_b$ of
such $p_a$ and $p_b$ span $(\Pol^{\rm f})^{H_I}$.  The polynomial $\tilde
q_I$ is $H_I$-invariant.  Let $y\in\Pol^{\rm f}$.  Then
\[
 \Av_H(\tilde q_Iy)=\frac{|H_I|}{|H|}\sum_{h\in H/H_I}h\big(\tilde
 q_I\,\Av_{H_I}(y)\big),
\]
and $\Av_{H_I}(y)$ is a combination of products $p_ap_b$.
So~\eqref{eq:cutclass} gives $[\mathsf S_{\lambda_m}(\Av_H(\tilde
q_Iy))]\in F_{l(m)-1}$.  The group $H$ acts on $\Pol^{\rm f}$ through $H_1$,
and $\epsilon^{\rm f}$ identifies $H_1$ with $\cW_m$, so
$\epsilon(\Av_H(\tilde q_Iy))=\Av_{\cW_m}(q_I\,\epsilon(y))$.  As $y$ varies,
$\epsilon(y)$ runs over $\Pol_m$.  The group $\cW_m$ permutes the edges of
$\Gamma_m$, so $J_m$ is $\cW_m$-stable, and by~\ref{I:alg}(a) these elements
span $\Av_{\cW_m}(J_m)=J_m^{\cW_m}$.
\end{details}

Now let $p\in\Pol^H$ with $\epsilon(p)\in J_m$.  By~(2), $\epsilon(p)\in
\Pol_m^{\cW_m}$, so $\epsilon(p)\in J_m^{\cW_m}$, and $\epsilon(p)=\epsilon(p')$
for a combination $p'$ of the $\Av_H(\tilde q_Iy)$.  By~(2), $[\mathsf
S_{\lambda_m}(p)]=[\mathsf S_{\lambda_m}(p')]\in F_{l(m)-1}$.

Finally, by~(2), $\epsilon(\Pol^H)=\Pol_m^{\cW_m}$, and
$(\Pol_m/J_m)^{\cW_m}=\Pol_m^{\cW_m}/J_m^{\cW_m}$ by~\ref{I:alg}(a).
So~(2) and~(3) give a well-defined map
$(\Pol_m/J_m)^{\cW_m}\langle\delta_m\rangle\to F_L/F_{L-1}$ for each $m$ with $l(m)=L$,
of degree zero.  The sum of these maps is onto, since $F_L$ is spanned by
$F_{L-1}$ and the $[\mathsf S_{\lambda_m}(p)]$ with $l(m)=L$.
\end{proof}

This is the argument of \pre{Prop.~5.2}.

\subsection{The induced filtration}\label{ss:induced}
Put $F'_L=\Phi(F_L)\subseteq H^*(X)$.  As $\Phi$ is onto, these subspaces
filter $H^*(X)$, with $F'_{k_0-1}=0$ and $F'_{|\bv|}=H^*(X)$, and $\Phi$
induces surjections $F_L/F_{L-1}\to F'_L/F'_{L-1}$.

\begin{lemma}\label{lem:reduction}
Suppose that for every $L$ the composite
\begin{equation}\label{eq:grcomposite}
 \bigoplus_{l(m)=L}(\Pol_m/J_m)^{\cW_m}\langle\delta_m\rangle
 \twoheadrightarrow F_L/F_{L-1}\twoheadrightarrow F'_L/F'_{L-1}
\end{equation}
of the surjection of Proposition~\ref{prop:filt} with the map induced by
$\Phi$ is injective.  Then Theorem~\ref{thm:loops} holds.
\end{lemma}

\begin{proof}
Both maps in~\eqref{eq:grcomposite} are then isomorphisms, so
$\dim B=\sum_L\dim F_L/F_{L-1}=\sum_L\dim F'_L/F'_{L-1}=\dim H^*(X)$, and the
surjection $\Phi$ is an isomorphism.  As $\Phi$ preserves degrees, the
Poincar\'e polynomial of $X$ is $\sum_mq^{\delta_m}h_m(q)$.
\end{proof}

The first case is $L=k_0$.  The only multipartition with $k_0$ parts is
$m_G$, with one part at each vertex, for which $\cW_{m_G}$ is trivial and
$\epsilon$ maps $\Pol^{\cW_{\lambda_G}}$ onto $\Pol_G$.  As $F_{k_0-1}=0$,
Proposition~\ref{prop:filt}(2) and~(3) give a well-defined map
\[
 \psi\colon(\Pol_G/J_G)\langle\delta_{m_G}\rangle\to H^*(X),\qquad
 \psi(\epsilon(p))=\Phi[\mathsf S_{\lambda_G}(p)]\quad\text{for }p\in\Pol^{\cW_{\lambda_G}},
\]
and this is the composite~\eqref{eq:grcomposite} for $L=k_0$.  The other
values of $L$ reduce to this case for other quivers, by restriction to Levi
subgroups (Proposition~\ref{prop:detect}), and in \S\ref{ss:top} we show that
$\psi$ is injective (Corollary~\ref{cor:topclass}).

\subsection{Levi subgroups}\label{ss:levi}
A Levi subgroup $M\supseteq T$ is given by a decomposition of each
$\Omega_k$ into blocks.  The gauge theory of $M$ acting on $\mathfrak
m\oplus N'$ is the quiver gauge theory of the following quiver $Q_M$:
\begin{itemize}[leftmargin=2em]
\item its vertices are the blocks, with dimension the size of the block;
\item the block $(k,D)$ has framing $W_k$ and one loop, the $(D,D)$ block of
$\ell_k$;
\item for each arrow $k\to l$ of $Q'$ and all blocks $(k,D)$ and $(l,E)$,
there is an arrow $(k,D)\to(l,E)$.
\end{itemize}
Keeping only the block-diagonal part $\mathfrak m$ of the chosen loops, rather
than all of $\mathfrak g$, gives $Q_M$ a loop at every vertex, so $Q_M$
satisfies the hypotheses of this section; the off-diagonal part is recovered
in Lemma~\ref{lem:ind}.  We write $A(M)$,
$B(M)$, $X(M)$, $\Phi_M$ and $\kappa_M$ for its Coulomb branch, B-algebra,
quiver variety at determinant stability, comparison map and Kirwan map.
Its Weyl group $\cW_M\subseteq\cW$ is the product of the symmetric groups
of the blocks.  We write $\mathsf S^M_\lambda(p)$ for its dressed monopoles,
which are sums over $\cW_M/\cW_{M,\lambda}$, where
$\cW_{M,\lambda}=\cW_M\cap\cW_\lambda$.  For $M=M_m$, built from the blocks used to
define $\lambda_m$, the multipartition of $Q_M$ with one part per block has
coweight $\lambda_m$, graph $\Gamma_m$, and data $\Pol_m$ and $J_m$.

\begin{lemma}[restriction]\label{lem:restr}
\begin{enumerate}[leftmargin=2em]
\item Proposition~\ref{prop:teleman} for $G$ and for $M$ gives $A\subseteq
A(M)$, with the same $r_\lambda$.  Followed by projection to the part of
$M$-charge zero, this induces a ring homomorphism
$\res_M\colon B\to B(M)$ with $\res_M[f]=[f]$ for $f\in R$.
\item If $l(m)$ is less than the number of blocks of $M$, then
$\res_M[\mathsf S_{\lambda_m}(p)]=0$.
\item $\res_{M_m}[\mathsf S_{\lambda_m}(p)]=[\mathsf
S^{M_m}_{\lambda_m}(p)]$, and this depends only on $\epsilon(p)$.  If
$m'\neq m$ and $l(m')=l(m)$, then $\res_{M_m}[\mathsf S_{\lambda_{m'}}(p)]=0$.
\end{enumerate}
\end{lemma}

\begin{proof}
(1) The quiver $Q_M$ has one chosen loop at each block, and removing these
loops leaves the $M$-module $N'$: the arrows of $Q_M$ between blocks are the
block components of the arrows of $Q'$, and its framings are the block
components of $\Hom(W,V)$.  So Proposition~\ref{prop:teleman} for $Q_M$ uses
the same torus $T$, the same weights $\cB'$ and the same $r_\lambda$, and it
gives $A(M)=\big(\bigoplus_\lambda\Pol\,r_\lambda\big)^{\cW_M}$.  Hence
$A\subseteq A(M)$.

The $M$-charge of $r_\lambda$ is the vector of the sums of $\lambda$ over the
blocks.  For $a\in A$, let $a_0$ be its component of $M$-charge zero.  Then
$\rho(a)=[a_0]\in B(M)$ is a ring homomorphism, and it kills the elements of
nonzero $G$-charge, since the $M$-charge refines the $G$-charge.  So it kills
$I_B$ and induces $\res_M\colon B\to B(M)$.  An element $f\in R$ has charge
zero, so $\res_M[f]=[f]$.
\begin{details}
As $\cW_M$ preserves each block, the $M$-charge components of an
element of $A(M)$ lie in $A(M)$, and $B(M)$ is the quotient of $A(M)$ by the
ideal generated by the components of nonzero $M$-charge.  If $f=\sum_cf_c$
and $g=\sum_cg_c$ are the decompositions by $M$-charge, then
$(fg)_0=f_0g_0+\sum_{c\neq0}f_cg_{-c}$, and each $f_c$ with $c\neq0$ vanishes
in $B(M)$.  The $G$-charge at $k$ is the
sum of the $M$-charges of the blocks at $k$, so an element of $A$ of nonzero
$G$-charge has only nonzero $M$-charge components.
\end{details}

For (2) and~(3), let $\lambda\in\Z^\Omega$ and $p\in\Pol^{\cW_\lambda}$.  Split
$\cW/\cW_\lambda$ into $\cW_M$-orbits.  The stabilizer in $\cW_M$ of
$w\cW_\lambda$ is $\cW_{M,w\lambda}=\cW_M\cap\cW_{w\lambda}$, and $wp$ is
$\cW_{w\lambda}$-invariant, so
\[
 \mathsf S_\lambda(p)=\sum_{w\in\cW_M\backslash\cW/\cW_\lambda}\mathsf
 S^M_{w\lambda}(wp).
\]
The term of $w$ has $M$-charge equal to the block sums of $w\lambda$, so
$\res_M$ keeps exactly the terms in which all block sums of $w\lambda$
vanish.

Take $\lambda=\lambda_m=s-\one$.  Then $w\lambda_m=ws-\one$, and $ws$ is $s_D$
at $w(d_D)$ for each part $D$ and $0$ elsewhere.  So the sum of $w\lambda_m$
over a block $\beta\subseteq\Omega_k$ is
$\sum_{D:\,w(d_D)\in\beta}s_D-|\beta|$.  For this to vanish, the parts $D$ with
$w(d_D)\in\beta$ must have sizes adding up to $|\beta|\ge1$, so at least one
$w(d_D)$ must lie in $\beta$.  If $l(m)$ is less than the number of blocks,
then for every $w$ some block contains no $w(d_D)$.  This proves~(2).

For~(3), let $M=M_m$, which has $l(m)$ blocks, the block of $D$ having size
$s_D$.  Let $l(m')=l(m)$.  If $w\lambda_{m'}$ has $M$-charge zero, then each
block contains exactly one of the elements $w(d_{D'})$, for $D'$ a part of
$m'$, and that part has the size of the block.  So at each vertex the parts of $m'$ are
the block sizes, which are the parts of $m$, and $m'=m$.  This proves the
last claim.  For $m'=m$, the condition says that on each block $w\lambda_m$
has one entry equal to the size of the block minus one and all other
entries $-1$.  These coweights form the $\cW_{M}$-orbit of $\lambda_m$, so only
the double coset of $w=1$ survives, and
$\res_{M_m}[\mathsf S_{\lambda_m}(p)]=[\mathsf S^{M_m}_{\lambda_m}(p)]$.  In
$Q_{M_m}$, $\lambda_m$ is the coweight of the multipartition with one part per
block, whose data are $\Zm$, $\Pol_m$ and $\epsilon$.  So
Proposition~\ref{prop:filt}(2) for $Q_{M_m}$ shows that this class depends
only on $\epsilon(p)$.
\end{proof}

\begin{lemma}[induction]\label{lem:ind}
Let $X(M)^\circ\subseteq X(M)$ be the open subset where, at each vertex $k$,
the blocks of $\ell_k$ have pairwise disjoint spectra.
\begin{enumerate}[leftmargin=2em]
\item There is a morphism $\iota_M\colon X(M)^\circ\to X$ with
$\iota_M^*\cV_k\cong\bigoplus_D\cV_{k,D}$, where $\cV_{k,D}$ is the tautological
bundle of the vertex $(k,D)$ of $Q_M$.  So
$\iota_M^*\kappa(f)=\kappa_M(f)|_{X(M)^\circ}$ for $f\in R$.
\item Restriction $H^*(X(M))\to H^*(X(M)^\circ)$ is injective.
\item $\iota_M^*\circ\Phi=\Phi_M\circ\res_M$, as maps from $B$ to
$H^*(X(M))\subseteq H^*(X(M)^\circ)$.
\end{enumerate}
\end{lemma}

\begin{proof}
Write a point of $T^*N=T^*(\mathfrak g\oplus N')$ as $(\ell,\ell^*,n)$, with
$\ell,\ell^*\in\mathfrak g$ the components for the chosen loops and $n\in
T^*N'$, and a point of $T^*(\mathfrak m\oplus N')$ as $(\ell_M,\ell_M^*,n)$
with $\ell_M,\ell_M^*\in\mathfrak m$.  Identify $\mathfrak g^*$ with
$\mathfrak g$ by the trace form, let $\mu'\colon T^*N'\to\mathfrak g$ be the
moment map of $G$ on $T^*N'$, and let $\mathrm{pr}_{\mathfrak m}$ and
$\mathrm{pr}_{\mathfrak m^\perp}$ be the block-diagonal and off-block-diagonal
parts.  Then
\[
 \mu(\ell,\ell^*,n)=[\ell,\ell^*]+\mu'(n),\qquad
 \mu_M(\ell_M,\ell_M^*,n)=[\ell_M,\ell_M^*]+\mathrm{pr}_{\mathfrak m}\mu'(n).
\]
Suppose the blocks of $\ell_M$ at each vertex have pairwise disjoint spectra.
For blocks $D\neq E$ at $k$ and $\gamma\in\mathfrak m^\perp$, the $(E,D)$ block of
$[\ell_M,\gamma]$ is $\ell_{k,E}\gamma_{ED}-\gamma_{ED}\ell_{k,D}$.  So by~\ref{I:sylv},
$\mathrm{ad}_{\ell_M}$ is invertible on $\mathfrak m^\perp$.  Put
\[
 \tilde\iota(\ell_M,\ell_M^*,n)=\big(\ell_M,\ \ell_M^*+\gamma,\ n\big),\qquad
 \gamma=-\mathrm{ad}_{\ell_M}^{-1}\mathrm{pr}_{\mathfrak m^\perp}\mu'(n)\in\mathfrak
 m^\perp.
\]
Since $[\ell_M,\ell_M^*]\in\mathfrak m$ and $[\ell_M,\gamma]\in\mathfrak m^\perp$,
\[
 \mu\big(\tilde\iota(\ell_M,\ell_M^*,n)\big)=[\ell_M,\ell_M^*]
 +\mathrm{pr}_{\mathfrak m}\mu'(n)+\big([\ell_M,\gamma]+\mathrm{pr}_{\mathfrak
 m^\perp}\mu'(n)\big)=\mu_M(\ell_M,\ell_M^*,n).
\]
So $\tilde\iota$ maps $\mu_M^{-1}(0)$ into $\mu^{-1}(0)$, and it is
$M$-equivariant.

It also preserves stability.  By~\ref{I:cyc}, a point of $\mu^{-1}(0)$ is
stable if and only if the subspace $U\subseteq V$ generated by the images of
the $W_k$ under the arrows of the doubled quiver of $Q$ is all of $V$.
Likewise for $Q_M$, with the subspace $U_M=\bigoplus_DU_M^{(D)}$ generated by
the images of the framings of the blocks under the arrows of the doubled
quiver of $Q_M$.  Let $U$ be the subspace for
$\tilde\iota(\ell_M,\ell_M^*,n)$.  Then $U=U_M$, because the projections onto
the blocks are polynomials in $\ell_M$, and by~\ref{I:sylv} the blocks of
$\gamma$ are combinations of compositions of block components of the arrows
of $Q'$, their duals and the framing maps.
\begin{details}
\begin{itemize}[leftmargin=2em]
\item $U\supseteq U_M$.  The subspace $U$ is stable under $\ell_k$, which is
$\ell_M$ at $k$.  By~\ref{I:sylv} the projections onto the $V^{(D)}$ are
polynomials in $\ell_k$, so $U=\bigoplus_D\mathrm{pr}_DU$.  Hence $U$ contains
the images $\mathrm{pr}_D(\mathrm{im}\,i_k)$ of the framings of the blocks,
and it is stable under the block components $\mathrm{pr}_E\circ
x\circ\mathrm{pr}_D$ of every arrow $x$ of the doubled quiver of $Q$.  These
block components include all arrows of the doubled quiver of $Q_M$.  For
the chosen loops, the diagonal blocks of $\ell_k$ and $\ell_k^*$ are
$\ell_{k,D}$ and $\ell^*_{k,D}$, as $\gamma$ has no diagonal blocks.
\item $U\subseteq U_M$.  We show that $U_M$ is stable under every arrow of
the doubled quiver of $Q$.  The arrows of $Q'$ and their duals, the framing
maps, $\ell_k$ and the block-diagonal part $\ell_M^*$ of $\ell_k^*$ are sums
of their block components, which preserve $U_M$.  It remains to treat $\gamma$.
By~\ref{I:sylv}, its $(E,D)$ block is a linear combination of maps
$\ell_{k,E}^r\circ\mathrm{pr}_E\,\mu'(n)\,\mathrm{pr}_D\circ\ell_{k,D}^s$.
The component of $\mu'(n)$ at $k$ is a sum of terms $\pm xx^*$ for arrows $x$
of $Q'$ and their duals, and of $i_kj_k$.  So
$\mathrm{pr}_E\,\mu'(n)\,\mathrm{pr}_D$ is a sum of compositions of block
components of arrows of $Q'$ and their duals, which preserve $U_M$, and of
$\mathrm{pr}_Ei_kj_k\mathrm{pr}_D$, whose image lies in
$\mathrm{pr}_E(\mathrm{im}\,i_k)\subseteq U_M$.  So $\gamma$ maps $U_M$ into
$U_M$.  Since $U_M$ contains the image of $W$, it contains $U$.
\end{itemize}
\end{details}
So $\tilde\iota(\ell_M,\ell_M^*,n)$ is stable if and only if
$(\ell_M,\ell_M^*,n)$ is.

Hence $\tilde\iota$ induces a morphism $\iota_M$ from $X(M)^\circ$, the
quotient by $M$ of the stable points of $\mu_M^{-1}(0)$ with disjoint block
spectra, to $X=\mu^{-1}(0)^{\rm st}/G$.  The tautological bundle $\cV_k$ on $X$
comes from the $G$-module $V_k$, which as an $M$-module is
$\bigoplus_DV^{(D)}$.  So $\iota_M^*\cV_k\cong\bigoplus_D\cV_{k,D}$.  The class
$\kappa(f)$ is $f$ evaluated at the Chern roots of the $\cV_k$, and $f\in
R\subseteq\Pol^{\cW_M}$, so $\iota_M^*\kappa(f)=\kappa_M(f)|_{X(M)^\circ}$.  This
gives~(1).

For~(2), let $\zeta\in\mathfrak z(\mathfrak m)$ be central, with entry
$\zeta_D$ on the block $D$.  The translation $\tau_\zeta\colon\ell_M\mapsto\ell_M+\zeta$
commutes with $M$, preserves $\mu_M$, since $[\zeta,\ell_M^*]=0$, and
preserves stability, since it does not change which subspaces are stable
under each $\ell_{k,D}$.  So $\tau_\zeta$ acts on $X(M)$, and it is homotopic to
the identity through the $\tau_{t\zeta}$, $t\in[0,1]$.  Let $\fL_M$ be the core
of $X(M)$.  By~\ref{I:core}(b), each block $\ell_{k,D}$ is nilpotent on
$\fL_M$.  Choose $\zeta$ with distinct entries at each vertex.  Then on
$\tau_\zeta(\fL_M)$ the block $\ell_{k,D}$ has the single eigenvalue $\zeta_D$,
so $\tau_\zeta(\fL_M)\subseteq X(M)^\circ$.  The composite
\[
 H^*(X(M))\to H^*(X(M)^\circ)\to H^*(\tau_\zeta(\fL_M))
 \xrightarrow{\ \tau_\zeta^*\ }H^*(\fL_M)
\]
is the pullback along $\tau_\zeta\colon\fL_M\to X(M)$, which equals restriction
to $\fL_M$ because $\tau_\zeta$ is homotopic to the identity.  That restriction is
an isomorphism by~\ref{I:core}(a).  So the first map is injective.

For~(3), every class of $B$ is $[f]$ for some $f\in R$, since $R\to B$ is
onto, and $\Phi[f]=\kappa(f)$.  By Lemma~\ref{lem:restr}(1),
$\res_M[f]=[f]$, and $\Phi_M[f]=\kappa_M(f)$.  So by~(1),
$\iota_M^*\Phi[f]=\kappa_M(f)|_{X(M)^\circ}=\Phi_M(\res_M[f])|_{X(M)^\circ}$,
and by~(2) restriction to $X(M)^\circ$ is injective.
\end{proof}

For $M=M_m$, the multipartition of $Q_{M_m}$ with one part per block plays the
role of $m_G$: its coweight, data and shift are $\lambda_m$, $\Pol_m$, $J_m$
and $\delta_m$, and the map $\psi$ of \S\ref{ss:induced} for $Q_{M_m}$ is
\[
 \psi_{M_m}\colon(\Pol_m/J_m)\langle\delta_m\rangle\to H^*(X(M_m)),\qquad
 \psi_{M_m}(\epsilon(p))=\Phi_{M_m}[\mathsf S^{M_m}_{\lambda_m}(p)].
\]

\begin{proposition}[detecting the pieces]\label{prop:detect}
Fix $L$, and for each multipartition $m$ with $l(m)=L$ take $M_m\supseteq T$
built from the blocks used to define $\lambda_m$.  The map
\[
 \Psi_L\colon F_L\to\bigoplus_{l(m)=L}H^*(X(M_m)),\qquad
 b\mapsto\big(\Phi_{M_m}(\res_{M_m}b)\big)_m=\big(\iota_{M_m}^*\Phi(b)\big)_m ,
\]
kills $F_{L-1}$ and factors through $F'_L/F'_{L-1}$.  Its composite with the
surjection of Proposition~\ref{prop:filt} is the direct sum over $m$ of the
restrictions of the $\psi_{M_m}$ to the $\cW_m$-invariants.  So if every
$\psi_{M_m}$ is injective, then so is~\eqref{eq:grcomposite}.
\end{proposition}

\begin{proof}
The two expressions for $\Psi_L$ agree by Lemma~\ref{lem:ind}(3), and the
second shows that $\Psi_L$ factors through $\Phi$.  Each $M_m$ has $L$ blocks,
so by Lemma~\ref{lem:restr}(2), $\Psi_L$ kills $F_{L-1}$; hence it factors
through $F'_L/F'_{L-1}$.  By Lemma~\ref{lem:restr}(3), $\Psi_L$ sends $[\mathsf
S_{\lambda_m}(p)]$ to $\Phi_{M_m}[\mathsf S^{M_m}_{\lambda_m}(p)]=\psi_{M_m}(\epsilon(p))$
in the summand of $m$, and to $0$ in the others.  A direct sum of injective
maps into different summands is injective.
\end{proof}

\subsection{The socle and the top degree}\label{ss:top}
By Proposition~\ref{prop:detect}, it remains to show that the map $\psi$ of
\S\ref{ss:induced} is injective, for every quiver with a loop at every vertex.
Let $\fL\subseteq X$ be the core, as in~\ref{I:core}, so that
$H^*(X)\cong H^*(\fL)$.  Put $d_{\max}=\dim N'-z-k_0$, the bound
$\delta_m+n_m-l(m)$ of Proposition~\ref{prop:toric} for $m=m_G$.

\begin{lemma}[socle]\label{lem:socle}
The map $\psi$ is $R$-linear, where $R$ acts on $\Pol_G/J_G$ through
$\epsilon$, which is onto, and on $H^*(X)$ through $\kappa$.  So $\psi$ is
injective if it is injective in the top degree of $\Pol_G/J_G$, which it maps
to degree $d_{\max}$.
\end{lemma}

\begin{proof}
For $f\in R$ and $p\in\Pol^{\cW_{\lambda_G}}$ we have $\mathsf
S_{\lambda_G}(fp)=f\,\mathsf S_{\lambda_G}(p)$, as $f$ is $\cW$-invariant, so
$\psi(\epsilon(f)\,\epsilon(p))=\kappa(f)\,\psi(\epsilon(p))$.  The map
$\epsilon\colon R\to\Pol_G$ is onto, since it sends $\sum_{d\in\Omega_k}x_d$ to
$v_kx_k$.  So $\ker\psi$ is an ideal of $\Pol_G/J_G$.  By~\ref{I:alg}(c), if
it is nonzero it meets the socle, which by Proposition~\ref{prop:toric} is the
top degree $n_{m_G}-k_0$ of $\Pol_G/J_G$.  The shift is $\delta_{m_G}$, and
$\delta_{m_G}+n_{m_G}-k_0=d_{\max}$ by Proposition~\ref{prop:toric}.
\end{proof}

\begin{proposition}\label{prop:top}
$B$ vanishes in degrees above $d_{\max}$.  In degree $d_{\max}$ it is
spanned by the classes $[\mathsf S_{\lambda_G}(p)]$, so $\psi$ maps the top
degree of $\Pol_G/J_G$ onto $H^{2d_{\max}}(X)$.  Also $\dim\fL\le d_{\max}$.
\end{proposition}

\begin{proof}
By Proposition~\ref{prop:filt}, $F_L/F_{L-1}$ is a quotient of
$\bigoplus_{l(m)=L}(\Pol_m/J_m)^{\cW_m}\langle\delta_m\rangle$.  By
Proposition~\ref{prop:toric}, $\Pol_m/J_m$ vanishes above degree
$n_m-l(m)$, so the summand of $m$ vanishes above degree
$\delta_m+n_m-l(m)=\dim N'-z-L$.  For $L=k_0$ this is $d_{\max}$, and the only
multipartition with $k_0$ parts is $m_G$, with one part at each vertex.  For
$L>k_0$ it is less than $d_{\max}$.  Since $F_{k_0-1}=0$ and $F_{|\bv|}=B$,
$B$ vanishes above degree $d_{\max}$, and its part of degree $d_{\max}$ is
that of $F_{k_0}$, which is spanned by the $[\mathsf S_{\lambda_G}(p)]$.  As
$\Phi$ is onto, $\psi$ maps onto $H^{2d_{\max}}(X)$.  If
$e=\dim\fL$ exceeded $d_{\max}$, then $H^{2e}(X)\cong H^{2e}(\fL)$ would be
nonzero by~\ref{I:core}(d).  This is degree $e$ in our grading, so the
surjectivity of $\Phi$ would contradict $B_e=0$.
\end{proof}

\begin{lemma}\label{lem:comp}
Let $o$ be an acyclic orientation of $\Gamma_{m_G}$, whose vertices are
$Q_0^+$ and $\infty$, in which $\infty$ is the only source.  Then $\fL$ contains an irreducible constructible subset
$\fL_o$, of dimension at least $d_{\max}$, with the following property.  For
every edge $k\to l$ of $o$ with $k,l\neq\infty$, the arrows of the doubled
quiver from $V_l$ to $V_k$ vanish on $\fL_o$, and at every point of $\fL_o$
some arrow from $V_k$ to $V_l$ is nonzero.  In particular
$\fL_o\not\subseteq\overline{\fL_{o'}}$ for $o\neq o'$.
\end{lemma}

\begin{proof}
In this proof $k$ and $l$ run over $Q_0^+$.  The parallel edges between two
vertices are oriented the same way, since $o$ is acyclic.  Choose a total order $\prec$ on $Q_0^+$ with $l\prec k$
whenever the edges between $k$ and $l$ point from $k$ to $l$, and bases
$e_{k,1},\dots,e_{k,v_k}$ of the $V_k$.  Order all the basis vectors so that
$e_{l,s}$ comes before $e_{k,r}$ if $l\prec k$, and $e_{k,s}$ before $e_{k,r}$
if $s<r$.  Let $V_{\le r}$ be the span of the first $r$ basis vectors, and let
$B_o=\prod_kB_k\subseteq G$ be the stabilizer of this flag, where $B_k$ is the
group of upper triangular matrices in the basis of $V_k$.

Let $\mathbb A\subseteq T^*N$ be the linear subspace where the maps $W\to V$
are arbitrary, the maps $V\to W$ vanish, and every arrow of the doubled
quiver maps $V_{\le r}$ into $V_{\le r-1}$ for all $r$.  So:
\begin{itemize}[leftmargin=2em]
\item the loops of the doubled quiver at $k$, namely $\ell_k$, $\ell_k^*$ and
the $c_{kk}$ loops coming from $Q'$, are strictly upper triangular;
\item for $l\prec k$, the $c_{kl}$ arrows of the doubled quiver from $V_k$
to $V_l$ are arbitrary, and the $c_{kl}$ arrows from $V_l$ to $V_k$ vanish.
\end{itemize}
Hence
\[
 \dim\mathbb A=\sum_k\Big((2+c_{kk})\tbinom{v_k}2+w_kv_k\Big)
 +\sum_av_{t(a)}v_{h(a)},
\]
the last sum over the arrows $a$ of $Q'$ that are not loops, and
$\dim B_o=\sum_k\binom{v_k+1}2$.  On $\mathbb A$, the component of $\mu$ at $k$
is strictly upper triangular with zeros just above the diagonal, so $\mu$ maps
$\mathbb A$ into a space of dimension $\sum_k\binom{v_k-1}2$.
\begin{details}
On
$\mathbb A$, the moment map at $k$ is a sum of terms $\pm xx^*$ over the arrows
$x$ at $k$, and $i_kj_k=0$.  For an arrow between $k$ and $l\neq k$, one of
$x$ and $x^*$ vanishes, so $xx^*=0$.  The remaining terms are commutators of
strictly upper triangular matrices.
\end{details}

Let $p_0\in\mathbb A$ be the following point:
\begin{itemize}[leftmargin=2em]
\item $\ell_k$ is the regular nilpotent $e_{k,r}\mapsto e_{k,r-1}$, and the
other loops are zero;
\item if the edges between $k$ and $l$ point from $k$ to $l$, then one
arrow of the doubled quiver from $k$ to $l$ is $e_{k,v_k}\mapsto
e_{l,v_l}$; all other components of arrows between different vertices are zero;
\item if $w_k>0$, some vector of $W_k$ maps to $e_{k,v_k}$, and the other
framing maps are zero.
\end{itemize}
Then $\mu(p_0)=0$, since $[\ell_k,\ell_k^*]=0$ and all other terms vanish.
Since $\infty$ is the only source, every $k$ has an incoming edge, either
from $\infty$, which means $w_k>0$, or from some $l\succ k$.  So by descending
induction along $\prec$, each $e_{k,v_k}$ lies in the subspace generated by
$W$, and hence so does $V_k=\C[\ell_k]e_{k,v_k}$.  So $p_0$ is stable.

Let $Y$ be an irreducible component of $\mathbb A\cap\mu^{-1}(0)$ through
$p_0$.  It is cut out of $\mathbb A$ by $\sum_k\binom{v_k-1}2$ equations, so
$\dim Y\ge\dim\mathbb A-\sum_k\binom{v_k-1}2$.  Let $Y^\circ\subseteq Y$ be the
subset of stable points at which every $\ell_k$ is regular nilpotent and,
for every pair $k,l$ whose edges point from $k$ to $l$, some component from
$V_k$ to $V_l$ is nonzero.  These conditions are open; on $\mathbb A$,
$\ell_k$ is nilpotent, and it is regular if and only if it has rank $v_k-1$.
They hold at $p_0$, so $Y^\circ$ is open and dense in $Y$.  The group $B_o$
preserves $\mathbb A$ and $\mu^{-1}(0)$, and it is connected, so it preserves
each irreducible component of $\mathbb A\cap\mu^{-1}(0)$, and it preserves
$Y^\circ$.

Let $\fL_o\subseteq X$ be the image of $Y^\circ$, which is irreducible and
constructible.  Let $\lambda$ be the cocharacter of $G$ acting on the $r$-th
basis vector by $t^{|\bv|+1-r}$.  Then $\lambda(t)\cdot p\to0$ as $t\to0$ for
every $p\in\mathbb A$, and $\fL_o\subseteq\fL$ by~\ref{I:core}(c).
\begin{details}
On $\mathbb A$, a component of an arrow from
the $r$-th to the $r'$-th basis vector, with $r'<r$, is scaled by
$t^{r-r'}$, and a framing map to the $r$-th basis vector by $t^{|\bv|+1-r}$.
\end{details}

If $p$ and $gp$ both lie in $Y^\circ$, then $g\in B_o$, since $g_k$ must
preserve the kernels of the powers of the regular nilpotent $\ell_k$.
\begin{details}
A strictly
upper triangular regular nilpotent $\ell$ on $V_k$ has
$\ker\ell^r=\mathrm{span}(e_{k,1},\dots,e_{k,r})$, since this span lies in the
kernel and the kernel has dimension $r$.  The loop of $gp$ at $k$ is
$g_k\ell_kg_k^{-1}$, so $g_k$ preserves these spans, and $g_k\in B_k$.
\end{details}
Stable points have trivial stabilizer, so the fibers of $Y^\circ\to\fL_o$ are
free $B_o$-orbits, and
\[
 \dim\fL_o=\dim Y^\circ-\dim B_o
 \ge\dim\mathbb A-\textstyle\sum_k\binom{v_k-1}2-\dim B_o=d_{\max}.
\]
\begin{details}
Since $2\binom v2-\binom{v-1}2-\binom{v+1}2=-1$, the right side is
\[
 \sum_k\big(\tfrac12c_{kk}v_k(v_k-1)+w_kv_k-1\big)+\sum_av_{t(a)}v_{h(a)}.
\]
This is $d_{\max}=\dim N'-z-k_0$: the loops of $Q'$ at $k$ contribute
$\frac12c_{kk}v_k^2$ to $\dim N'$, and $z=\sum_k\frac12c_{kk}v_k$.
\end{details}

For the last claim, let $o\neq o'$.  All edges at $\infty$ point away from
$\infty$ in both, so there are $k,l\in Q_0^+$ whose edges point from $k$ to
$l$ in $o$ and from $l$ to $k$ in $o'$.  The condition that all arrows of the
doubled quiver from $V_k$ to $V_l$ vanish is closed and $G$-invariant.  It
holds on $\fL_{o'}$, hence on its closure, and it fails at every point of
$\fL_o$.
\end{proof}

\begin{corollary}\label{cor:topclass}
The map $\psi$ of \S\ref{ss:induced} is injective.  If $\Gamma_{m_G}$ has an
acyclic orientation in which $\infty$ is the only source, then
$\dim\fL=d_{\max}$, and the closures of the $\fL_o$ are distinct irreducible
components of $\fL$.
\end{corollary}

\begin{proof}
If $\Gamma_{m_G}$ has no orientation $o$ as above, then $\Pol_G/J_G$ vanishes
in its top degree $n_{m_G}-k_0$, by Proposition~\ref{prop:toric}.  If it were
nonzero, its socle would be nonzero by~\ref{I:alg}(c) and would lie in that
degree.  So $\Pol_G/J_G=0$, and there is nothing to prove.

Otherwise, Lemma~\ref{lem:comp} and Proposition~\ref{prop:top} give
$\dim\fL=d_{\max}$.  So each $\overline{\fL_o}$ is an irreducible closed
subset of $\fL$ of dimension $\dim\fL$, hence a component, and they are
distinct by Lemma~\ref{lem:comp}.  By~\ref{I:core}(d), $\dim
H^{2d_{\max}}(X)$ is at least the number of orientations $o$, which by
Proposition~\ref{prop:toric} is the dimension of the top degree of
$\Pol_G/J_G$.  By Proposition~\ref{prop:top}, $\psi$ maps that top degree onto
$H^{2d_{\max}}(X)$, so it is injective there, and by Lemma~\ref{lem:socle} it
is injective.
\end{proof}

\subsection{Proof of Theorem~\ref{thm:loops}}\label{ss:proof}
Fix $L$.  For each multipartition $m$ with $l(m)=L$,
Corollary~\ref{cor:topclass} for the quiver $Q_{M_m}$ shows that $\psi_{M_m}$
is injective.  By Proposition~\ref{prop:detect}, the
composite~\eqref{eq:grcomposite} is injective, and
Lemma~\ref{lem:reduction} gives the theorem.\qed

\medskip
The proof also shows that the surjection of Proposition~\ref{prop:filt} is an
isomorphism
$F_L/F_{L-1}\cong\bigoplus_{l(m)=L}(\Pol_m/J_m)^{\cW_m}\langle\delta_m\rangle$, and
that $\Psi_L$ is injective on $F_L/F_{L-1}$.

\begin{corollary}\label{cor:cusp}
Let $\mathcal M_L$ be the set of Levi subgroups $M\supseteq T$ with at
least $L$ blocks.  Then $F_{L-1}=\bigcap_{M\in\mathcal M_L}\ker\res_M$.  In
particular the \emph{cuspidal part}
$\bigcap_{M\neq G}\ker\res_M$ of $B$ is $F_{k_0}\cong(\Pol_G/J_G)\langle\delta_{m_G}\rangle$.
This is the shifted cohomology of the hypertoric variety of the center of
$G$ acting on $N'^{\xi_{m_G}}$.
\end{corollary}

\begin{proof}
The inclusion $\subseteq$ is Lemma~\ref{lem:restr}(2), since $F_{L-1}$ is
spanned by classes $[\mathsf S_{\lambda_m}(p)]$ with $l(m)<L$.  Conversely, let
$b$ lie in the right side, and let $L'$ be minimal with $b\in F_{L'}$.
Suppose $L'\ge L$.  Each $M_m$ with $l(m)=L'$ has $L'\ge L$ blocks, so
$\Psi_{L'}(b)=0$, and as $\Psi_{L'}$ is injective on $F_{L'}/F_{L'-1}$, we get
$b\in F_{L'-1}$, a contradiction.  So $L'\le L-1$.  For the cuspidal part,
the Levi subgroups $M\supseteq T$ other than $G$ are those with at least
$k_0+1$ blocks, so the intersection is $F_{k_0}$.  It is isomorphic to
$(\Pol_G/J_G)\langle\delta_{m_G}\rangle$ by the proof of Theorem~\ref{thm:loops}, as
$F_{k_0-1}=0$, and the last claim is Proposition~\ref{prop:toric}.
\end{proof}

\begin{remark}[Hausel's formula]\label{rem:hausel}
For quivers with a loop at every vertex, the formula of
Theorem~\ref{thm:loops} for all $\bv$ at once is equivalent to Hausel's
formula \cite{hauselBettiNumbers2006}, given the identity of
Remark~\ref{rem:kirwan}(4).  \pre{\S6} uses this in the other direction.  We
follow \cite[\S3]{nakajimaMultipartitionSpaces2026}.

Fix $\bw$, and let $m$ run over all tuples $(m^{(k)})_{k\in Q_0}$ of
partitions, of any sizes, with $|m|=(|m^{(k)}|)_k$.  For partitions put
$\langle\mu,\nu\rangle=\sum_{a,b}\min(\mu_a,\nu_b)$, and put
\[
 d_m=\sum_{a\in Q'_1}\langle m^{(t(a))},m^{(h(a))}\rangle ,
\]
which is $\dim N_{Q'}^{\xi_m}$ for
$N_{Q'}=\bigoplus_{a\in Q'_1}\Hom(V_{t(a)},V_{h(a)})$ and $\bv=|m|$, by the
computation before Definition~\ref{def:toric-data}.  So $\dim
N'^{\xi_m}=d_m+\bw\cdot l(m)$, where $l(m)$ is now the vector of the numbers of
parts at each vertex.  Consider the series in variables $z_k$, $k\in Q_0$,
\[
 A_{\rm fr}=\sum_m z^{|m|}\frac{q^{-d_m-\bw\cdot l(m)}}{(q)_m},\qquad
 A_0=\sum_m z^{|m|}\frac{q^{-d_m}}{(q)_m},\qquad
 H=\sum_m z^{|m|}q^{-d_m-\bw\cdot l(m)}h_m(q).
\]
\begin{itemize}[leftmargin=2em]
\item Let $M_\eta=\mu^{-1}(\eta)/\!/G$ be the affine quotient at
$\eta=(\mathrm{id}_{V_k})_k$, of dimension $2\delta(\bv,\bw)$.  For any quiver,
Hausel's formula is
\[
 \sum_{\bv}q^{-\delta(\bv,\bw)}P_{M_\eta}(q)\,z^{\bv}
 =\frac{\sum_m z^{|m|}q^{-E(m)-\bw\cdot l(m)}/(q)_m}{\sum_m z^{|m|}q^{-E(m)}/(q)_m},
\]
in the form of \pre{(6.9)}, where
$E(m)=\sum_{a\in Q_1}\langle m^{(t(a))},m^{(h(a))}\rangle-\sum_k\langle m^{(k)},m^{(k)}\rangle$.  For our $Q$, the chosen loops cancel the last
sum, so $E(m)=d_m$, and similarly $\delta(\bv,\bw)=\dim N'$.  By~\ref{I:rot}
with $\theta=\nu_{\det}$, followed by multiplying the cotangent coordinate by
$i$, $M_\eta$ is homeomorphic to $X$.  So Hausel's formula says
$\sum_{\bv}q^{-\dim N'}P_X(q)\,z^{\bv}=A_{\rm fr}/A_0$.
\item Theorem~\ref{thm:loops} says that this sum is $H$, since
$q^{-\dim N'}q^{\delta_m}=q^{-\dim N'^{\xi_m}}$.
\item Multiplying Remark~\ref{rem:kirwan}(4) by $z^{|m|}q^{-d_m-\bw\cdot l(m)}$
and summing over $m$ gives $A_{\rm fr}=A_0H$, as the term of a pair $(m,I)$
factors into a term of $A_0$ and a term of $H$.
\begin{details}
Let $m_I$ and $m'$ be
the multipartitions formed by the parts in $I$ and outside $I$, and let $e$
be the number of edges of $\Gamma_m$ between $I$ and the other parts.  Then
$d_m=d_{m_I}+d_{m'}+e$ and $\deg q_I=e+\bw\cdot l(m_I)$, so the term of $(m,I)$
is the product of the term of $m_I$ in $A_0$ and that of $m'$ in $H$.  Also
$(m,I)\mapsto(m_I,m')$ is a bijection from pairs of a multipartition and a
$\cW_m$-orbit of sets of its parts onto pairs of multipartitions.
\end{details}
\end{itemize}
Since $A_0$ has constant term $1$, the first two statements are equivalent,
given the third.

In Hausel's formula the multipartitions arise in a different way.  Hausel's
proof uses an arithmetic Fourier transform.  In the point count of an
unpublished 2004 note of Nakajima, explained in \cite[Rem.~3.2]{nakajimaMultipartitionSpaces2026}, one
counts isomorphism classes of framed and of unframed representations of $Q$
over a finite field by Burnside's lemma.  The sum over the elements $g$ of
$G(\mathbb F)$ factors over the eigenvalues of $g$, the framing changes only
the factor of the eigenvalue $1$, and the term of $m$ in $A_{\rm fr}$ or $A_0$
is the contribution of the unipotent $g$ of Jordan type $m$, after
substituting $|\mathbb F|^{-1}$ for $q$.  By the Krull--Schmidt theorem, the
ratio counts the absolutely indecomposable framed representations, and these
count the points of $M_\eta$ \cite{crawley-boeveyAbsolutelyIndecomposable2004}.  In
Theorem~\ref{thm:higgs} below, $m$ records instead the multiplicities of the
eigenvalues of the chosen loops.
\end{remark}

\section{Deformations and quantization}\label{sec:deform}
This section follows \pre{\S7}.  Write $H_F=H^*_F(\mathrm{pt})$ and
$H_{F,\hbar}=H^*_{F\times\C^\times_\hbar}(\mathrm{pt})$.

\begin{lemma}[graded Nakayama]\label{lem:flavor}
Let $\nu$ be generic, and suppose $\Phi\colon B^\nu_0\to H^*(X_\nu)$ is an
isomorphism.  Then for every flavor torus $F$, the maps
$\Phi\colon B^\nu_F\to H^*_F(X_\nu)$ and
$\Phi\colon B^\nu(A_{F,\hbar})\to H^*_{F\times\C^\times_\hbar}(X_\nu)$ are
isomorphisms.
\end{lemma}

\begin{proof}
We treat the second map; the first is the same with $H_F$ in place of
$H_{F,\hbar}$.  Let $H^+$ be the part of $H_{F,\hbar}$ of positive degree,
and let $K$ be the kernel of the surjection $\Phi$ of Theorem~\ref{thm:map}.
Tensor the exact sequence
$0\to K\to B^\nu(A_{F,\hbar})\to H^*_{F\times\C^\times_\hbar}(X_\nu)\to0$ with
$\C=H_{F,\hbar}/H^+$.  The last term is free over $H_{F,\hbar}$
by~\ref{I:kirwan}, so the sequence stays exact:
\[
 0\to K/H^+K\to B^\nu(A_{F,\hbar})/H^+\to H^*(X_\nu)\to0 .
\]
By~\ref{I:base}, the middle term is $B^\nu_0$, and the right map is the
isomorphism $\Phi$ for $F$ trivial and $\hbar=0$.  So $K=H^+K$.  As $K$ is a
graded submodule of a quotient of $R_\hbar$, it is bounded below, and graded
Nakayama gives $K=0$; here $R_\hbar\to B^\nu(A_{F,\hbar})$ is onto by
Lemma~\ref{lem:qsurj}.
\end{proof}

\begin{lemma}\label{lem:betti}
The Betti numbers of $X_\theta$ do not depend on the generic parameter
$\theta$.
\end{lemma}

\begin{proof}
Let $Z\cong\C^{Q_0}$ be the space of central values of the moment map, and
let $\pi\colon\mu^{-1}(Z)/\!/G\to Z$ be the family of affine quotients.  As
$R\pi_!\Q$ is constructible, there is a Zariski open dense
$U\subseteq Z$ over which the compactly supported Betti numbers of the fibers
$\mu^{-1}(\zeta)/\!/G$ are constant; $U$ is connected.  Let $\theta$ be
generic.  Every $\theta'$ in a small real ball around $\theta$ is generic,
with the same signs of the $\theta'\cdot\bb$, so $X_{\theta'}=X_\theta$ by
\S\ref{ss:setup}.  By~\ref{I:rot}, $X_{\theta'}$ is homeomorphic to
$\mu^{-1}(-i\theta')/\!/G$.  The points $-i\theta'$ form a real ball in
$i\mathbb R^{Q_0}$, which is Zariski dense in $Z$, so some $-i\theta'$ lies in
$U$.  Hence the compactly supported Betti numbers of $X_\theta$ are those of
the fibers over $U$, for every generic $\theta$.  Poincar\'e duality on the
smooth $X_\theta$ gives the ordinary Betti numbers.
\end{proof}

The $X_\theta$ for generic $\theta$ are even diffeomorphic
\cite[Cor.~4.2]{nakajimaInstantonsALE1994}, by the same hyperk\"ahler rotation.  That paper
assumes that the quiver has no loops, but, as \cite[Rem.~3.3]{nakajimaMultipartitionSpaces2026} points out, the proof of
\cite[Cor.~4.2]{nakajimaInstantonsALE1994} uses only
\cite[Thm.~4.1(1),(2)]{nakajimaInstantonsALE1994} and the rotation of complex
structures, and these do not need that assumption.

\emph{Abelianization with root denominators.}  When $Q$ does not have a loop
at every vertex, Proposition~\ref{prop:teleman} is not available, and we use
the abelianization map of \cite{BFN} directly.  Use $T$, $\Omega$,
$\Omega_k$, $\Pol$ and $\cW$ as in \S\ref{ss:abel}, let $F$ be a flavor
torus, and put $\Pol_F=H^*_{T\times F}(\mathrm{pt})$.  Let $\cB$ be the
multiset of $T\times F$-weights of $N$, and define
$r_\lambda\in\mathrm{Frac}(\Pol_F)[\Z^\Omega]$ as in \S\ref{ss:abel}, with $\cB$
in place of $\cB'$; here $\langle\beta,\lambda\rangle$ pairs the $T$-part of
$\beta$ with $\lambda$.  For the rest of this section $r_\lambda$ and $\mathsf
S_\lambda$ have the meanings given here.  These $r_\lambda$ multiply by~\eqref{eq:rr} with $\cB$
in place of $\cB'$, so their $\Pol_F$-span is the abelian Coulomb branch of
$(T,N)$ \cite[Thm.~4.1]{BFN}.  The abelianization map $(\iota_*)^{-1}$ is an
injective ring homomorphism from $A_F$ to $\mathrm{Frac}(\Pol_F)[\Z^\Omega]$
\cite[Lemma~5.17]{BFN}, and the charge corresponds to the charge of
$\lambda$.

Call $\lambda\in\Z^\Omega$ \emph{minuscule} if on each $\Omega_k$ its entries
take at most two values, differing by one, and put
\[
 \mathrm{eu}_\lambda=\prod_{k\in Q_0}\ \prod_{d,e\in\Omega_k,\ \lambda_d>\lambda_e}(x_d-x_e).
\]
For minuscule $\lambda$ and $p\in\Pol_F^{\cW_\lambda}$ put
\[
 \mathsf S_\lambda(p)=\sum_{w\in\cW/\cW_\lambda}w\Big(\frac{p}{\mathrm{eu}_\lambda}\,r_\lambda\Big).
\]
For such $\lambda$ the orbit $\Gr^\lambda$ is closed, its $T$-fixed points
are the $w\lambda$, and at $\hbar=0$ the Euler class of its tangent space at
$w\lambda$ is $\pm w(\mathrm{eu}_\lambda)$.  So by \cite[Prop.~6.6]{BFN},
$(\iota_*)^{-1}$ sends the dressed monopole $M_{\lambda,p}$ of~\ref{I:gen} to
$\pm\mathsf S_\lambda(p)$, with a sign depending only on $\lambda$.  In the
setting of \S\ref{ss:abel}, the weights of the chosen loops contribute
$\pm\mathrm{eu}_\lambda$ to $r_\lambda$ for minuscule $\lambda$, so these $\mathsf
S_\lambda(p)$ agree up to sign with those defined there.

\begin{proposition}\label{prop:loopremoval}
For every quiver, $\Phi\colon B_0\to H^*(X_{\nu_{\det}})$ is an
isomorphism.
\end{proposition}

\begin{proof}
We follow \pre{Prop.~7.3}.  Write $X=X_{\nu_{\det}}$.  The maps $R\to B_0$
and $\kappa\colon R\to H^*(X)$ are onto, by Lemma~\ref{lem:generic}
and~\ref{I:kirwan}, and $\Phi[f]=\kappa(f)$.  So we must show that their
kernels $K_B$ and $K_X$ agree.

Let $Q^+$ be $Q$ with a loop $\ell_k^+$ added at every vertex.  Let
$\C^\times_u$ scale $\ell_k^+$ with weight $1$ and its dual $\ell_k^{+*}$ with
weight $-1$.  This commutes with $G$ and preserves $\mu$, since
$[t\ell_k^+,t^{-1}\ell_k^{+*}]=[\ell_k^+,\ell_k^{+*}]$, so $\C^\times_u$ is a
flavor torus.  Write $H^*_u$ for $\C^\times_u$-equivariant cohomology and $u$ for
the equivariant parameter, so that $H^*_{G\times\C^\times_u}(\mathrm{pt})=R[u]$.
Let $A^+$ and $B^+$ be the Coulomb branch of $Q^+$ with flavor $\C^\times_u$
at $\hbar=0$ and its B-algebra, and let $X^+=X_{\nu_{\det}}(Q^+)$, with Kirwan
map $\kappa^+\colon R[u]\to H^*_u(X^+)$.  Every vertex of $Q^+$ carries a loop,
so by Theorem~\ref{thm:loops} and Lemmas~\ref{lem:flavor}
and~\ref{lem:generic}, the kernel $K^+$ of $R[u]\to B^+$ is the kernel of
$\kappa^+$.  Let $\mathfrak n\subseteq R[u]$ be the maximal ideal generated by
the positive part $R_+$ of $R$ and by $u-1$.  We show that $K_X[u]$, $K^+$ and
$K_B[u]$ have the same localization at $\mathfrak n$.

\emph{Quiver variety side.}  Let $p$ be a stable point of the zero fiber of
the moment map of $Q^+$ whose class is fixed by $\C^\times_u$.  As $G$ acts
freely on stable points, $t\cdot p=\rho(t)\cdot p$ for a unique $\rho(t)\in G$,
and $\rho$ is a cocharacter of $G$.  Its conjugacy class is locally constant,
so the fixed locus $(X^+)^{\C^\times_u}$ is the disjoint union of the open and
closed subsets $Y_\rho$ on which it is the class of $\rho$.  Let $n_d$,
$d\in\Omega$, be the weights of $\C^\times_u$ on the fibers of the tautological
bundles over $Y_\rho$.  If $\rho\neq1$, then $f=\sum_{d\in\Omega}x_d^2$ and
$u^2\sum_dn_d^2$ act on $H^*_u(Y_\rho)$ by elements that differ by a nilpotent,
and $f-u^2\sum_dn_d^2\in R[u]$ does not lie in $\mathfrak n$, as it is
$-\sum_dn_d^2\neq0$ modulo $\mathfrak n$.  So $H^*_u(Y_\rho)_{\mathfrak n}=0$.
\begin{details}
$\C^\times_u$ acts on the fibers of $\cV_k$ through $\rho^{-1}$.  As
$\C^\times_u$ acts trivially on $Y_\rho$,
$H^*_u(Y_\rho)=H^*(Y_\rho)[u]$, and the equivariant Chern roots of $\cV_k$ there
are $y_d+n_du$, with $y_d$ the ordinary Chern roots.  So for $f\in R$,
$\kappa^+(f)|_{Y_\rho}-f\big((n_du)_d\big)$ lies in $H^{>0}(Y_\rho)[u]$, which is
a nilpotent ideal.
\end{details}
If $\rho=1$, then $t\cdot p=p$, so
$\ell^+=\ell^{+*}=0$.  For such $p$ the moment map and the stability condition
of~\ref{I:cyc} are those of $Q$.  So $Y_1=X$, with $\C^\times_u$ acting
trivially on $X$ and on its tautological bundles, and $\kappa^+$ followed by
restriction to $Y_1$ is $\kappa$ extended linearly over $\C[u]$.

By the localization theorem \cite[Thm.~3]{brionEquivariantCohomology1998}, restriction
$H^*_u(X^+)\to H^*_u\big((X^+)^{\C^\times_u}\big)$ becomes an isomorphism after
inverting $u$.
\begin{details}
Its hypothesis, an equivariant embedding into a
finite-dimensional representation of the compact torus $U(1)\subseteq\C^\times_u$,
holds because $X^+$ is a $\C^\times_u$-stable closed subvariety of
$\mathbb P(U)\times U'$ for some $\C^\times_u$-modules $U$ and $U'$, by its
construction as a GIT quotient.
\end{details}
As $u$ is a unit at $\mathfrak n$,
\[
 H^*_u(X^+)_{\mathfrak n}\cong H^*_u(Y_1)_{\mathfrak n}=(H^*(X)[u])_{\mathfrak n},
\]
compatibly with $\kappa^+$ and $\kappa$.  Localization is exact, so
$K^+_{\mathfrak n}=K_X[u]_{\mathfrak n}$.

\emph{Coulomb side.}  Use the abelianization above for $Q$ without flavor,
extended linearly over $\C[u]$, and for $Q^+$ with flavor $\C^\times_u$.  It
embeds $A[u]$ and $A^+$ in $\mathrm{Frac}(\Pol[u])[\Z^\Omega]$, compatibly with
the charge.  The weights of $Q^+$ are those of $Q$ together with the weights
$x_e-x_d+u$ of the added loops, for $d,e$ in the same $\Omega_k$, and the root
factors $\mathrm{eu}_\lambda$ do not change.  So $r^+_\lambda=g_\lambda r_\lambda$,
where
\[
 g_\lambda=\prod_{k\in Q_0}\ \prod_{d,e\in\Omega_k}(x_e-x_d+u)^{(\lambda_d-\lambda_e)_+}.
\]
The factor $g_\lambda$ is $\cW_\lambda$-invariant, so for minuscule $\lambda$
the elements $\mathsf S^+_\lambda(p)$ of $Q^+$, defined with $r^+_\lambda$ in
place of $r_\lambda$, are $\mathsf S^+_\lambda(p)=\mathsf S_\lambda(g_\lambda p)$.
By~\ref{I:gen} these generate $A^+$ over $R[u]$, so $A^+\subseteq A[u]$.
Conversely, $\sigma_\lambda=\prod_{w\in\cW/\cW_\lambda}w(g_\lambda)$ lies in $R[u]$
and takes the value $1$ at $x=0$, $u=1$, so $\sigma_\lambda\notin\mathfrak n$.
The quotient $h_\lambda=\sigma_\lambda/g_\lambda$ is a $\cW_\lambda$-invariant
polynomial, and as $\sigma_\lambda$ is $\cW$-invariant,
\[
 \mathsf S_\lambda(p)=\sigma_\lambda^{-1}\,\mathsf S_\lambda(g_\lambda h_\lambda p)
 =\sigma_\lambda^{-1}\,\mathsf S^+_\lambda(h_\lambda p).
\]
By~\ref{I:gen} the $\mathsf S_\lambda(p)$ generate $A[u]$ over $R[u]$, so
$A^+_{\mathfrak n}=A[u]_{\mathfrak n}$ inside $\mathrm{Frac}(\Pol[u])[\Z^\Omega]$.
The B-algebra is the quotient by the ideal generated by the elements of
nonzero charge, which commutes with localization.  So $B^+_{\mathfrak
n}=(B_0[u])_{\mathfrak n}$, compatibly with the maps from $R[u]_{\mathfrak n}$,
whose images are the multiples of $e^0$ in both.  Hence $K^+_{\mathfrak
n}=K_B[u]_{\mathfrak n}$.

\emph{Conclusion.}  We have shown $K_B[u]_{\mathfrak n}=K_X[u]_{\mathfrak n}$,
and graded Nakayama gives $K_B=K_X$.
\begin{details}
Let $E=(K_B+K_X)/K_B$, a finitely generated graded $R$-module.  Then
$E\otimes_RR[u]_{\mathfrak n}=0$.  Tensoring further with $R[u]_{\mathfrak
n}/\mathfrak n=\C$, on which $R$ acts through $R/R_+$, gives $E/R_+E=0$, so
$E=0$ by graded Nakayama.  So $K_X\subseteq K_B$, and by symmetry $K_B=K_X$.
\end{details}
\end{proof}

\begin{theorem}\label{thm:classical}
Let $\nu$ be generic.  Then $B^\nu_F=B_F$, and the maps
$\Phi\colon B_F\to H^*_F(X_\nu)$ and
$\Phi\colon B^\nu(A_{F,\hbar})\to H^*_{F\times\C^\times_\hbar}(X_\nu)$ are
isomorphisms.  So the maps $R\to B^\nu_F$ and $\kappa\colon R\to H^*_F(X_\nu)$
are surjective with the same kernel, and Theorem~\ref{thm:main} holds.
\end{theorem}

\begin{proof}
By Lemma~\ref{lem:generic}, $B^\nu_F=B_F$ and $B^\nu_0=B_0$, which does not
depend on $\nu$.  By Proposition~\ref{prop:loopremoval} and
Lemma~\ref{lem:betti}, $\dim B_0=\dim H^*(X_{\nu_{\det}})=\dim H^*(X_\nu)$.  So
the surjection $\Phi\colon B^\nu_0\to H^*(X_\nu)$ of \S\ref{sec:map} is an
isomorphism, and Lemma~\ref{lem:flavor} gives the rest.  For the kernels,
the map $R\to B^\nu_F$ is onto by
Lemma~\ref{lem:generic}, and $\Phi[f]=\kappa(f)$, so its kernel is that of
$\kappa$.
\end{proof}

% Input from Hikita-rewrite.tex.
\section{A second proof of the dimension bound}\label{sec:bps}

% COMMENT: drafted 2026-10-14 from two unattributed notes Ben received,
% hikita_bps_only_proof.pdf and hall_coulomb_comparison.pdf (not in the
% repository).  Credit is still to be settled.  Citation checks against the
% sources are recorded in NOTES.md, "The BPS route".

This section gives a second proof of Proposition~\ref{prop:loopremoval}:
$\dim B_0=\dim H^*(X_{\nu_{\det}})$ for every quiver.  It adds no loops and
does not use \S\ref{sec:loops}.  Instead, it writes $B_0$ as a quotient of a
space $P_\Bd$: the part of dimension $(\bv,1)$ of a cohomological Hall algebra
of the framed quiver $\hat Q$, modulo decomposable elements and the image of
one tautological operator.  BPS cohomology spans $P_\Bd$, and Botta and
Davison identify it with $H^*(X_{\nu_{\det}})$.  The argument gives the
dimension only, not the Poincar\'e polynomial of Theorem~\ref{thm:loops}.

The proof is conditional.  It uses results of Jindal and Neguț
\cite{jindalLoopNilpotent2026} whose proofs are only sketched there; see
Remark~\ref{rem:bps-status}.  Throughout, $F$ is trivial, $\hbar=0$ on the
Coulomb side, $X=X_{\nu_{\det}}$, and $B_0=A_0/I_B$, where $I_B$ is generated
by the elements of nonzero charge.

\subsection{The framed quiver}\label{ss:bps-framed}
Let $\hat Q$ be $Q$ with one more vertex $\infty$ and $w_i$ arrows from
$\infty$ to $i$, as in \S\ref{ss:setup}, and put $\Bd=(\bv,1)$.  Let
$V_\infty=\C$ and $\hat G=G\times\GL(V_\infty)$.  After choosing bases of the
$W_i$, the representation space of $\hat Q$ of dimension $\Bd$ is $N$, with
$\Hom(W_i,V_i)=\Hom(V_\infty,V_i)^{w_i}$.  Let $\hat A$ be the Coulomb branch
of $(\hat G,N)$, at $\hbar=0$ and without flavor.  Its charges lie in
$\Z^{Q_0}\times\Z$.  The scalars in $\hat G$ act trivially on $N$, and this
splits $\hat A$, as in \pre{(8.4)}.

We use the abelianization of \S\ref{sec:deform}, as set out before
Proposition~\ref{prop:loopremoval}, for $\hat Q$ and $\Bd$ as well as for $Q$
and $\bv$.  Put $\hat\Omega=\Omega\sqcup\{\infty\}$, and let $\hat x_d$,
$d\in\hat\Omega$, be the coordinates on the Lie algebra of the diagonal torus
$\hat T$ of $\hat G$.  Put
\[
 x_d=\hat x_d-\hat x_\infty\quad(d\in\Omega),\qquad x_\infty=\hat x_\infty,
\]
so that $H^*_{\hat T}(\mathrm{pt})=\Pol[x_\infty]$.  Let $\hat\one\in\Z^{\hat\Omega}$
be the coweight with all entries $1$, and put
\[
 b=\sum_{d\in\Omega}x_d\in R,\qquad D=1+\sum_{k\in Q_0}v_k .
\]
Write $\hat r_\lambda$ and $\hat{\mathsf S}_\lambda(p)$ for the elements of
\S\ref{sec:deform} defined for $(\hat T,N)$, and $r_\lambda$, $\mathsf
S_\lambda(p)$ for those of $(T,N)$, written in the variables $x_d$.

\begin{lemma}\label{lem:bps-split}
Identify $e^{(\mu,0)}$ with $e^\mu$ for $\mu\in\Z^\Omega$.  Then:
\begin{enumerate}[leftmargin=2em]
\item $\hat r_{(\mu,k)}=r_{\mu-k\one}\,e^{k\hat\one}$ for $(\mu,k)\in\Z^\Omega\times\Z$;
in particular $\hat r_{\pm\hat\one}=e^{\pm\hat\one}$;
\item $\hat A=A_0[x_\infty,e^{\pm\hat\one}]$, and the map
$A_0\otimes\C[x_\infty,e^{\pm\hat\one}]\to\hat A$ is an isomorphism;
\item the part of $\hat A$ of charge $(\gamma,k)$ is
$A_0^{(\gamma-k\bv)}[x_\infty]\,e^{k\hat\one}$;
\item $\sum_{d\in\hat\Omega}\hat x_d=b+Dx_\infty$.
\end{enumerate}
\end{lemma}

\begin{proof}
The $\hat T$-weights of $N$ are the $\hat x_e-\hat x_d=x_e-x_d$, for the arrows
of $Q$, and the $\hat x_d-\hat x_\infty=x_d$, for the framing.  In the
variables $x_d$ these are the $T$-weights $\cB$ of $N$, and each pairs with
$(\mu,k)$ as it pairs with $\mu-k\one$.  This gives~(1), as
$e^{(\mu,k)}=e^{(\mu-k\one,0)}e^{k\hat\one}$.  Since $V_\infty$ is a line,
the root factors also agree: $\mathrm{eu}_{(\mu,k)}=\mathrm{eu}_{\mu-k\one}$.
So $\hat{\mathsf S}_{(\mu,k)}(p)=\mathsf S_{\mu-k\one}(p)\,e^{k\hat\one}$ for
$p\in\Pol^{\cW_\mu}[x_\infty]$.

A coweight $(\mu,k)$ has all entries in $\{0,1\}$ exactly when
$k\in\{0,1\}$ and $\mu-k\one$ has all entries in $\{0,1\}$ (if $k=0$) or in
$\{-1,0\}$ (if $k=1$); similarly for $\{0,-1\}$.  Also
$H^*_{\hat G}(\mathrm{pt})=R[x_\infty]$.  So~\ref{I:gen}, for $\hat Q$ and for $Q$,
shows that $\hat A$ is generated over $R[x_\infty]$ by the
$\mathsf S_\nu(p)\,e^{k\hat\one}$, with $\nu$ minuscule and $k\in\{0,\pm1\}$,
and that these generate $A_0[x_\infty,e^{\pm\hat\one}]$.  This is the first
part of~(2).  The map in~(2) is injective because
$\Z^{\hat\Omega}=\Z^\Omega\times\Z\hat\one$ and $x_\infty$ is transcendental
over $\mathrm{Frac}(\Pol)$.  For~(3), the coweight $(\mu,k)$ has charge
$(c,k)$, where $c$ is the charge of $\mu$, and $\mu-k\one$ has charge
$c-k\bv$.  For~(4), $\sum_{d\in\Omega}(x_d+x_\infty)+x_\infty=b+Dx_\infty$.
\end{proof}

Define an algebra homomorphism
\[
 \Theta\colon\hat A\to B_0,\qquad a\,x_\infty^je^{k\hat\one}\mapsto[a]\,(-b/D)^j\qquad(a\in A_0).
\]
It is well defined by Lemma~\ref{lem:bps-split}(2).  It is an algebra map
because the isomorphism there is one of algebras, and $\Theta$ is the tensor
product of the quotient map $A_0\to B_0$ with
$x_\infty\mapsto-b/D$, $e^{\hat\one}\mapsto1$.
Lemma~\ref{lem:bps-split}(3),(4) give:
\begin{enumerate}[leftmargin=2.6em,label=($\Theta$\arabic*)]
\item\label{Th:charge} $\Theta$ kills the part of $\hat A$ of charge
$(\gamma,k)$ unless $\gamma=k\bv$;
\item\label{Th:det} $\Theta\big((\sum_{d\in\hat\Omega}\hat x_d)\,a\big)=0$ for all
$a\in\hat A$;
\item\label{Th:coef} $\Theta\big(p(\hat x_d-\hat x_\infty:d\in\Omega)\big)=[p]$ for
$p\in R$.
\end{enumerate}
The value $-b/D$ is the only one compatible with~\ref{Th:det}.  The
corresponding projection $\hat x\mapsto\hat x-\frac1D(\sum_d\hat x_d)\hat\one$
splits off the scalars in the Lie algebra of $\hat T$.  Since it divides by
$D$, it does not come from a splitting of $\hat T$, which is why we use
$\Theta$ and not a decomposition of $\hat G$.

\begin{lemma}\label{lem:bps-charge}
If $0\neq\alpha\le\Bd$ and $\alpha\neq\Bd$, write $\alpha=(\gamma,k)$.  Then
$\gamma\neq k\bv$.
\end{lemma}

\begin{proof}
Here $k\in\{0,1\}$.  If $k=0$ then $\gamma\neq0$.  If $k=1$ then
$\gamma\le\bv$ and $\gamma\neq\bv$.
\end{proof}

\subsection{The loop-nilpotent Hall algebra}\label{ss:bps-hall}
Let $\mathsf P=\C[\hbar,u_a:a\in\hat Q_1]$, with maximal ideal
$\fm=(\hbar,u_a)$.  Here $\hbar$ is a new parameter on the Hall algebra side; it
becomes the loop rotation parameter of the Coulomb branch, which we then set
to $0$.  The tripled quiver of $\hat Q$ has the arrows $a$ of $\hat Q$, their
reverses $a^*$, and a loop $\omega_k$ at every vertex $k$.  It carries its
canonical potential and the action of a torus with weights $-u_a$ on $a$,
$u_a-\hbar$ on $a^*$ and $\hbar$ on $\omega_k$, which preserves the potential
\cite[\S3.1]{jindalLoopNilpotent2026}.  Let
$\cH_{\mathsf P}=\bigoplus_{\Bn}\cH_{\mathsf P,\Bn}$ be the loop-nilpotent
cohomological Hall algebra of \cite[\S3]{jindalLoopNilpotent2026}, with
coefficients extended to $\C$.  It is the equivariant critical cohomological
Hall algebra of the tripled quiver, with the $\omega_k$ required to be
nilpotent.  It is graded by dimension vectors
$\Bn\in\Z_{\ge0}^{\hat Q_0}$, with $\cH_{\mathsf P,0}=\mathsf P$, and each
$\cH_{\mathsf P,\Bn}$ is a module over
$H^*_{G_\Bn}(\mathrm{pt})\otimes\mathsf P=\mathsf P[z_{k,a}:a\le n_k]^{\mathrm{sym}}$, where
$G_\Bn=\prod_k\GL_{n_k}$.  We use the product $F\star G$, the product of
\cite{jindalLoopNilpotent2026} twisted by the sign
$(-1)^{\sum_{a\colon k\to l}n_lm_k}$ for $F$ of dimension $\Bn$ and $G$ of
dimension $\Bm$, as in \pre{(8.1)}.  The twist changes a word only by a sign
depending on the dimensions of its letters, so it does not change spans of
words.

For $g\in\mathsf P[z]^{\mathrm{sym}}$ let $e_\Bn(g)=g\cdot[\omega=0]_\Bn\in\cH_{\mathsf
P,\Bn}$ \cite[(98)]{jindalLoopNilpotent2026}.  Under the shuffle map of
\cite[\S3]{jindalLoopNilpotent2026} it goes to
$g(z)\prod_k\prod_{a,b=1}^{n_k}(z_{k,b}-z_{k,a}+\hbar)$
\cite[Lemma~3.6]{jindalLoopNilpotent2026}.  By definition,
$e_\Bn(hg)=h\cdot e_\Bn(g)$ for symmetric $h$.  Let $C_\Bn$ be the action of
$\sum_{k,a}z_{k,a}=c_1(\bigotimes_k\det V_k)$, the first Chern class of the
total determinant, on $\cH_{\mathsf P,\Bn}$.  It is a derivation:
\begin{equation}\label{eq:bps-derivation}
 C_{\Bn+\Bm}(F\star G)=(C_\Bn F)\star G+F\star(C_\Bm G),
\end{equation}
because the total determinant is multiplicative in short exact sequences;
\cite[Prop.~4.1]{davisonAffineBPS2022} proves this vertex by vertex.

Let $\cH=\cH_{\mathsf P}/\fm\cH_{\mathsf P}$, graded by the $\cH_\Bn$, and write
$e_\Bn(g)$ and $C_\Bn$ also for their images.  This specialization is
abstract.  The shuffle image of $e_\Bn(g)$ is divisible by $\hbar^{|\Bn|}$, so
it vanishes in the specialized shuffle algebra, but the class of $e_\Bn(g)$ in
$\cH$ need not vanish.  Put
\[
 \cT_\Bd=\sum_{\substack{\alpha+\beta=\Bd\\ \alpha,\beta\neq0}}\cH_\alpha\star\cH_\beta,
 \qquad
 P_\Bd=\cH_\Bd\big/(\cT_\Bd+C_\Bd\cH_\Bd).
\]
By~\eqref{eq:bps-derivation}, $C_\Bd\cT_\Bd\subseteq\cT_\Bd$.

We use the following inputs.
\begin{enumerate}[leftmargin=2.6em,label=(H\arabic*)]
\item\label{H:poly} \emph{Polynomial generation}
\cite[Thm.~2.4, Lemma~3.6, Prop.~3.7, Thm.~3.5]{jindalLoopNilpotent2026}.
$\cH_{\mathsf P}$ is generated as a $\mathsf P$-algebra by the elements
$e_\Bn(g)$.
\item\label{H:bps} \emph{BPS generation} \cite[(101)--(110),
Cor.~3.15]{jindalLoopNilpotent2026}.  There are $\mathsf P$-modules $L_\Bn$, the
loop-nilpotent BPS cohomology, and $\mathsf P$-linear maps
$\beta_\Bn\colon L_\Bn\to\cH_{\mathsf P,\Bn}$, such that $\cH_{\mathsf P}$ is
generated as a $\mathsf P$-algebra by the $C_\Bn^k\beta_\Bn(\ell)$, for $k\ge0$
and $\ell\in L_\Bn$.  This rests on cohomological integrality
\cite[Thm.~C]{davisonCohomologicalDonaldsonThomas2020}, in its equivariant
form \cite[Thm.~11.7]{davisonBPSAlgebras2023},
\cite[Thm.~3.4]{davisonAffineBPS2022}.  The tautological class $u$ of
\cite[(103)]{jindalLoopNilpotent2026} acts as $C_\Bn$; this is the
convention of \cite[\S3.4, (18)--(19)]{davisonAffineBPS2022}.
\item\label{H:framed} \emph{Framed BPS cohomology.}  If $\bw\neq0$, then
$L_\Bd$ is a free $\mathsf P$-module of rank $\dim H^*(X)$.  Indeed,
$L_\Bn\cong\hbar M_\Bn$ as $\mathsf P$-modules
\cite[Prop.~3.14]{jindalLoopNilpotent2026}, where $M_\Bn$ is the equivariant
BPS Lie algebra of the preprojective algebra of $\hat Q$ in dimension $\Bn$
\cite[Def.~3.6]{bottaOkounkovConjecture2023}.  Next,
$M_\Bn\cong M_\Bn^{\mathrm{noneq}}\otimes\mathsf P$
\cite[Cor.~11.10]{davisonBPSAlgebras2023}; the group allowed there contains
our torus.  Finally, $M_\Bd^{\mathrm{noneq}}\cong H^*(X)$, up to a shift of
degrees \cite[Cor.~3.13]{bottaOkounkovConjecture2023}.  The last result uses
framing arrows from $\infty$ to $i$, as here, and the stability $\delta_\infty$,
for which a module of dimension $\Bd$ is semistable exactly when no proper
submodule contains $V_\infty$.  So it is semistable exactly when the images of
the $W_i$ generate $V$, which is $\nu_{\det}$-stability by~\ref{I:cyc}.
\end{enumerate}

\begin{remark}[status of the inputs]\label{rem:bps-status}
With independent parameters, the genericity conditions of
\cite[\S2.2]{jindalLoopNilpotent2026} hold.
\begin{enumerate}[leftmargin=2em]
\item \cite[Thm.~2.4]{jindalLoopNilpotent2026} is a statement about a
shuffle algebra, and its proof is only sketched: it is ``completely
analogous'' to the K-theoretic case treated in earlier work of the same
authors.  The transfer to $\cH_{\mathsf P}$ uses
\cite[Prop.~3.7]{jindalLoopNilpotent2026} (proof sketched),
\cite[Lemma~3.6]{jindalLoopNilpotent2026} (stated without proof) and the
injectivity of the shuffle map \cite[Thm.~3.5]{jindalLoopNilpotent2026} (proof
omitted).
\item The PBW isomorphism for the loop-nilpotent algebra is
\cite[(109)]{jindalLoopNilpotent2026}, stated as an isomorphism of graded
vector spaces.  That it is induced by the product, which is what generation
needs, is implicit.  The proof of \cite[Cor.~3.15]{jindalLoopNilpotent2026}
cites the PBW theorem for the full algebra, Thm.~3.10 there, where (109) is
needed.
\item \cite[Prop.~3.14]{jindalLoopNilpotent2026} has a complete proof, which
uses \cite[Thm.~3.5]{jindalLoopNilpotent2026}.  The identification
\cite[Cor.~3.13]{bottaOkounkovConjecture2023} is sheaf-theoretic.  It combines
a comparison of BPS sheaves for different stability conditions
\cite[Prop.~3.4]{bottaOkounkovConjecture2023} with the description of the
stable framed BPS sheaf as the intersection complex of $X\times\A^1$
\cite[Prop.~3.11]{bottaOkounkovConjecture2023}, where $\A^1$ records the
common scalar of the loops $\omega_k$.  No Kac polynomials, point counts or
Hausel's formula enter.
\end{enumerate}
\end{remark}

\subsection{The Hall--Coulomb map}\label{ss:bps-map}
For $\Bn\le\Bd$ let $\lambda_\Bn\in\Z^{\hat\Omega}$ be the coweight that is $1$ on
the first $n_k$ elements of each $\hat\Omega_k$ and $0$ elsewhere.  It is
minuscule, with charge $\Bn$, and $\lambda_\Bd=\hat\one$.  For $g$ in
$\C[z_{k,a}:a\le n_k]^{\mathrm{sym}}$ write $g(\hat x_\Bn)$ for $g$
evaluated at the $\hat x_d$ with $(\lambda_\Bn)_d=1$.  It lies in
$\Pol[x_\infty]^{\cW_{\lambda_\Bn}}$.

\begin{lemma}[\pre{Lemma~8.1}]\label{lem:hall-coulomb}
Assume~\ref{H:poly}.  There is an algebra homomorphism
$\Upsilon\colon(\cH,\star)\to\hat A$ with the following properties.  It sends
$\cH_\Bn$ into the part of charge $\Bn$, and $e_\Bn(g)$ to $0$ if
$\Bn\not\le\Bd$.  For $\Bn\le\Bd$,
\[
 \Upsilon(e_\Bn(g))=\epsilon_\Bn\,\hat{\mathsf S}_{\lambda_\Bn}\big(g(\hat x_\Bn)\big),
\]
with $\epsilon_\Bn=\pm1$ depending only on $\Bn$.  In particular
$\Upsilon(e_\Bd(g))=\epsilon_\Bd\,g(\hat x)\,e^{\hat\one}$.
\end{lemma}

\begin{proof}[Sketch]
We follow \pre{Lemma~8.1}, which treats the shuffle algebra; it applies to
$\cH_{\mathsf P}$ through the shuffle map, an algebra homomorphism.  Over
$\mathrm{Frac}(\mathsf P)$, Tsymbaliuk's homomorphism from the shuffle
algebra to difference operators
\cite[Thm.~5.7]{tsymbaliukDifferenceOperators2023} gives an algebra
homomorphism $\Upsilon_{\mathsf P}$ from $(\cH_{\mathsf P},\star)$ to
difference operators in the $\hat x_d$, with
$D^\lambda f(\hat x)=f(\hat x+\hbar\lambda)D^\lambda$.  It sends $e_\Bn(g)$ to
$0$ if $\Bn\not\le\Bd$, and otherwise to $\pm E_\Bn(g)$, where
\[
 E_\Bn(g)=\sum_{w\in\cW/\cW_{\lambda_\Bn}}w\Big(g(\hat x_\Bn)\,
 \frac{\prod_{a\colon k\to l}\prod(\hat x_d-\hat x_e+u_a)}
      {\prod_k\prod(\hat x_d-\hat x_e)}\Big)D^{w\lambda_\Bn}
 \qquad\text{\pre{(8.2)}}.
\]
In the numerator $d\in\hat\Omega_k$ and $e\in\hat\Omega_l$; in the
denominator $d,e\in\hat\Omega_k$.  In both, $(\lambda_\Bn)_d=1$ and
$(\lambda_\Bn)_e=0$.  The reason is that Tsymbaliuk's map evaluates a shuffle
element on strings with step $\hbar$.  The factor $\prod(z_{k,b}-z_{k,a}+\hbar)$
vanishes on every string of length at least two, so only minuscule shifts
survive.  Its diagonal factors $\hbar^{|\Bn|}$ cancel the normalization of
\pre{Lemma~8.1}.

The coefficients of the $E_\Bn(g)$ lie in $\mathsf P[\hat x]$, localized at the
$\hat x_d-\hat x_e+r\hbar$ for $d\neq e$ in the same $\hat\Omega_k$ and
$r\in\Z$.  This ring is stable under the shifts, so the difference operators
with coefficients in it form a $\mathsf P$-algebra $\cD_{\mathsf P}$.  By~\ref{H:poly},
$\Upsilon_{\mathsf P}(\cH_{\mathsf P})\subseteq\cD_{\mathsf P}$, with no parameter
inverted.  Reducing modulo $\fm$ makes the $D^\lambda$ commute, and turns
$\cD_{\mathsf P}\otimes_{\mathsf P}\mathsf P/\fm$ into
$\Pol[x_\infty][(\hat x_d-\hat x_e)^{-1}][\Z^{\hat\Omega}]$, with
$D^\lambda\mapsto e^\lambda$.  Call the resulting map $\Upsilon$.

At $u_a=0$ the denominator of $E_\Bn(g)$ is $\mathrm{eu}_{\lambda_\Bn}$, and the
numerator is, up to a sign depending on $\Bn$, the product of the weights
$\beta$ of $N$ with $\langle\beta,\lambda_\Bn\rangle=-1$.  This is the
prefactor of $\hat r_{\lambda_\Bn}$.  So $\Upsilon(e_\Bn(g))$ is
$\pm\hat{\mathsf S}_{\lambda_\Bn}(g(\hat x_\Bn))$, which is the image of a dressed
minuscule monopole \cite[Prop.~6.6]{BFN}.  The images of the generators lie
in $\hat A$, and by~\ref{H:poly} so does the image of $\cH$.  The charge of
$w\lambda_\Bn$ is $\Bn$.  For $\Bn=\Bd$ both products are empty and
$\cW_{\lambda_\Bd}=\cW$, so $E_\Bd(g)=g(\hat x)D^{\hat\one}$.
\end{proof}

% COMMENT: [plausible, not checked against Tsymbaliuk's formula] Step 3
% below can avoid (H1).  At hbar = 0 the image of h F, for h symmetric,
% multiplies the D^lambda term of the image of F by h evaluated on the
% strings, so the image of C_d F is delta(image of F), where delta multiplies
% e^lambda by <lambda, \hat x>; that is, delta is minus the Poisson bracket
% with the Casimir (1/2) sum \hat x_d^2.  Then Theta(delta(a)) = 0 for a of
% charge d: delta(e^{\hat 1}) = (sum \hat x_d) e^{\hat 1}, and on A_0^{(0)} =
% R + sum A^{(-c)}A^{(c)} (Lemma lem:generic) delta kills R and preserves I_B.
% This is the derivation D_{i,k} of the preprint's section 5.  (H1) would
% then be needed only for the image of Upsilon to lie in \hat A.

\subsection{The bound}\label{ss:bps-bound}

\begin{theorem}\label{thm:bps}
Assume~\ref{H:poly}, \ref{H:bps} and~\ref{H:framed}.  For every quiver $Q$,
loops and multiple arrows allowed, and all $\bv$ and $\bw$,
\[
 \dim B_0\le\dim H^*(X).
\]
\end{theorem}

\begin{proof}
If $\bv=0$, then $B_0=\C$ and $X$ is a point.  If $\bw=0$ and $\bv\neq0$,
then $X$ is empty by~\ref{I:cyc}.  Also $\one$ pairs to zero with every weight
of $N$, so $r_{\pm\one}=e^{\pm\one}$, and by~\ref{I:gen} these lie in $A_0$.
They are inverse to each other and have nonzero charge $\pm\bv$, so $1\in I_B$
and $B_0=0$.  So assume $\bv\neq0$ and $\bw\neq0$, and let
\[
 \psi=\Theta\circ\Upsilon|_{\cH_\Bd}\colon\cH_\Bd\to B_0 .
\]

\emph{Step 1: $\psi$ is onto.}  For $p\in R$, let $g=p(z_{k,a}-z_{\infty,1})$,
which is symmetric at each vertex.  By Lemma~\ref{lem:hall-coulomb},
$\Upsilon(e_\Bd(g))=\epsilon_\Bd\,p(\hat x_d-\hat x_\infty)\,e^{\hat\one}$, so
$\psi(e_\Bd(g))=\epsilon_\Bd[p]$ by~\ref{Th:coef}.  The map $R\to B_0$ is onto
by Lemma~\ref{lem:generic}.

\emph{Step 2: $\psi(\cT_\Bd)=0$.}  Let $\alpha+\beta=\Bd$ with
$\alpha,\beta\neq0$.  Then $\Upsilon(\cH_\alpha)$ has charge $\alpha$, which
$\Theta$ kills by~\ref{Th:charge} and Lemma~\ref{lem:bps-charge}.  As $\Upsilon$
and $\Theta$ are multiplicative, $\psi(\cH_\alpha\star\cH_\beta)=0$.

\emph{Step 3: $\psi(C_\Bd\cH_\Bd)=0$.}  By~\ref{H:poly}, $\cH$ is spanned by
words in the $e_\Bn(g)$.  Letters with $\Bn=0$ are scalars, and a word of
dimension $\Bd$ with at least two letters of nonzero dimension lies in
$\cT_\Bd$.  So $\cH_\Bd$ is spanned by the $e_\Bd(g)$ and $\cT_\Bd$.  Now
$C_\Bd e_\Bd(g)=e_\Bd\big((\sum_{k,a}z_{k,a})g\big)$, so
\[
 \psi(C_\Bd e_\Bd(g))=\epsilon_\Bd\,\Theta\Big(\big(\textstyle\sum_{d\in\hat\Omega}\hat x_d\big)\,g(\hat x)\,e^{\hat\one}\Big)=0
\]
by~\ref{Th:det}.  Also $C_\Bd\cT_\Bd\subseteq\cT_\Bd$, which $\psi$ kills by
Step~2.  So $\psi$ induces a surjection $P_\Bd\to B_0$.

\emph{Step 4: the bound.}  Put $\bar L_\Bd=L_\Bd/\fm L_\Bd$.  By~\ref{H:bps},
and as tensor products are right exact, $\cH$ is generated by the images of
the $C_\Bn^k\beta_\Bn(\ell)$.  The word argument of Step~3 shows that $\cH_\Bd$
is spanned by $\cT_\Bd$ and the $C_\Bd^k\beta_\Bd(\ell)$, $\ell\in\bar L_\Bd$.
The terms with $k\ge1$ lie in $C_\Bd\cH_\Bd$, so $\bar L_\Bd\to P_\Bd$ is onto.
By~\ref{H:framed}, $\dim\bar L_\Bd=\dim H^*(X)$.  Hence
\[
 \dim B_0\le\dim P_\Bd\le\dim\bar L_\Bd=\dim H^*(X).\qedhere
\]
\end{proof}

No injectivity is needed anywhere in this proof, and no map is compared with
a geometric one.

\begin{corollary}\label{cor:bps}
Assume~\ref{H:poly}, \ref{H:bps} and~\ref{H:framed}.  Then
Proposition~\ref{prop:loopremoval} holds: for every quiver,
$\Phi\colon B_0\to H^*(X)$ is an isomorphism.  For $\bw\neq0$ the
surjections $\bar L_\Bd\to P_\Bd\to B_0$ of the proof are isomorphisms of
vector spaces.  With this in place of Proposition~\ref{prop:loopremoval}, the
proof of Theorem~\ref{thm:classical}, and so of Theorem~\ref{thm:main}, goes
through without \S\ref{sec:loops}.
\end{corollary}

\begin{proof}
The surjection $\Phi\colon B_0=B^{\nu_{\det}}_0\to H^*(X)$ of \S\ref{sec:map}
and Theorem~\ref{thm:bps} give $\dim B_0=\dim H^*(X)$, finite, so $\Phi$ is an
isomorphism.  Then every map in the chain $\bar L_\Bd\to P_\Bd\to B_0$ is a
surjection between spaces of the same finite dimension.
\end{proof}

\begin{remark}\label{rem:bps-compare}
Compared with the proof through \S\ref{sec:loops}:
\begin{enumerate}[leftmargin=2em]
\item This proof uses \S\ref{sec:map}, Lemmas~\ref{lem:flavor}
and~\ref{lem:betti}, and inputs~\ref{I:gen}--\ref{I:cyc}.  It does not use
\S\ref{sec:loops}, the loops added in Proposition~\ref{prop:loopremoval}, or
inputs~\ref{I:hypertoric}--\ref{I:core}.
\item It needs~\ref{H:poly}--\ref{H:framed} and \pre{Lemma~8.1}.  The proof in
the rest of the paper does not use Jindal--Neguț at all, because
Lemma~\ref{lem:generic} replaces \pre{\S8}.
\item It gives $\dim H^*(X)$, but not the formula of Theorem~\ref{thm:loops} or
the multipartition filtration.  It does not equip $\bar L_\Bd$ or $P_\Bd$ with
the ring structure of $B_0$, and we have not checked whether the maps
$\bar L_\Bd\to P_\Bd\to B_0$ respect gradings.
\end{enumerate}
\end{remark}


\section{The quiver variety side}\label{sec:higgs}
In the setting of \S\ref{sec:loops}, the filtration of
Proposition~\ref{prop:filt} and the formula of Theorem~\ref{thm:loops} can
be described on $X$ alone.  In this section cohomology carries its
cohomological grading, so the degrees of \S\ref{sec:loops} are doubled: $x_D$
has degree $2$, and $\langle\delta\rangle$ raises degrees by $2\delta$.  Let $\fc=\prod_{k\in Q_0^+}\A^{v_k}$ be the
space of tuples of monic polynomials of degrees $v_k$, and let
\[
 \chi\colon X\to\fc,\qquad [p]\mapsto\big(\det(t-\ell_k)\big)_k
\]
record the characteristic polynomials of the chosen loops.  The function
$\det(t-\ell_k)$ is $G$-invariant, so $\chi$ is well defined on $X$.  For the
scaling action on $X$, $\chi$ is equivariant for an action of $\C^\times$ on
$\fc$ with positive weights, contracting $\fc$ to the point $0=(t^{v_k})_k$.
\begin{details}
The scaling action of $a\in\C^\times$ on $T^*N$ multiplies $\ell_k$ by
$a$, and hence the coefficient of $t^{v_k-i}$ in $\det(t-\ell_k)$ by $a^i$.
\end{details}
The map $\chi$ is
not proper, so we prove the purity of $R\chi_*\Q_X$ by hand.

\emph{Strata.}  For a multipartition $m$, let
\[
 \pi_m\colon\A^{P_m}\to\fc,\qquad(c_D)_D\mapsto\Big(\prod_{D:\,k(D)=k}(t-c_D)^{s_D}\Big)_k ,
\]
and let $\fc_m$ be the image of the set where the $c_D$ at each vertex are
pairwise distinct.  So $\fc_m$ is the locus where the roots at each vertex
have multiplicities given by $m$.
\begin{itemize}[leftmargin=2em]
\item The map $\pi_m$ is finite, since the $c_D$ are roots of the image
polynomials, and it factors through $\A^{P_m}/\cW_m$, where $\cW_m$ permutes
the $c_D$ for parts of equal size at the same vertex.
\item On the set where the $c_D$ at each vertex are distinct, $\pi_m$ is a
Galois cover of $\fc_m$ with group $\cW_m$, the cover by labeled roots.  So
$\fc_m$ is smooth of dimension $l(m)$, and a $\cW_m$-module $V$ gives a local
system $\cL(V)$ on $\fc_m$.  For graded $V$ we write $\cL(V)$ and
$\IC(\overline{\fc_m},\cL(V))$ for the sums over the degrees $i$ of
$\cL(V^i)[-i]$ and $\IC(\overline{\fc_m},\cL(V^i))[-i]$.
\item The closure $\overline{\fc_m}=\pi_m(\A^{P_m})$ is the union of the
$\fc_{m'}$ over the coarsenings $m'$ of $m$, obtained by merging parts at the
same vertex, as happens when roots collide.  The map
$\A^{P_m}/\cW_m\to\overline{\fc_m}$ is finite and birational, and its source
is normal, so it is the normalization.  Only the point $c=0$ of
$\A^{P_m}/\cW_m$ lies over $0\in\fc$.
\end{itemize}
Let $\fc_{\ge L}=\bigcup_{l(m)\ge L}\fc_m$, the locus where the loops have at
least $L$ distinct eigenvalues in all, counting each vertex separately.  It is
open.

\emph{Levi subgroups.}  For a Levi subgroup $M\supseteq T$, let
$\chi_M\colon X(M)\to\fc_M=\prod_\beta\A^{|\beta|}$ record the characteristic
polynomials of the blocks $\ell_{k,\beta}$ of the chosen loops, the product
over the blocks $\beta$.  Let $\Delta_M\subseteq X(M)$ be the closed subset
where each block of each loop has a single eigenvalue.  The group $N_G(M)$
acts on the gauge theory of $M$, so $N_G(M)/M$ acts on $X(M)$, preserving
$\Delta_M$.  The \emph{cuspidal part} of $H^*(X(M))$ is the
$N_G(M)/M$-module
\[
 V_M=\ker\big(H^*(X(M))\to H^*(X(M)\setminus\Delta_M)\big).
\]
By Theorem~\ref{thm:higgs}(3) below, $\Phi_M$ identifies it with the cuspidal
part of $B(M)$ in the sense of Corollary~\ref{cor:cusp}.
As in Corollary~\ref{cor:cusp}, $\mathcal M_L$ is the set of Levi subgroups
$M\supseteq T$ with at least $L$ blocks.

\begin{theorem}\label{thm:higgs}
\begin{enumerate}[leftmargin=2em]
\item $R\chi_*\Q_X$ is pure of weight $0$, and
\[
 R\chi_*\Q_X\cong\bigoplus_m\IC\big(\overline{\fc_m},\cL(V_{M_m})\big)[-l(m)].
\]
\item $H^*(X)\cong\bigoplus_mV_{M_m}^{\cW_m}$.  For every $L$, the kernel
of restriction $H^*(X)\to H^*(\chi^{-1}\fc_{\ge L})$ is
$\bigcap_{M\in\mathcal M_L}\ker\iota_M^*=\Phi(F_{L-1})$.
\item $V_{M_m}=\Phi_{M_m}\big(F_{l(m)}(M_m)\big)\cong(\Pol_m/J_m)\langle\delta_m\rangle$, as
graded $\cW_m$-modules, where $F_\bullet(M_m)$ is the filtration of
Proposition~\ref{prop:filt} for $Q_{M_m}$.
\end{enumerate}
\end{theorem}

So the formula of Theorem~\ref{thm:loops} is the decomposition theorem for
$\chi$, read at the cone point.  The summand of $m$ lives on the closure of
the stratum where the eigenvalues of the loops have multiplicities $m$; this
records the centralizers of the semisimple parts of the loops, not their
Jordan types.  Its multiplicity space is the cohomology of the hypertoric variety of
$\Zm$ acting on $N'^{\xi_m}$, shifted by $\delta_m$, with its
$\cW_m$-monodromy, and the stalk at $0$ keeps the invariants.

We prove Theorem~\ref{thm:higgs} after two lemmas.  Write $K=R\chi_*\Q_X$ and
$K_M=R\chi_{M*}\Q_{X(M)}$, and use the conventions of~\ref{I:sheaves}.

\begin{lemma}[local model]\label{lem:localmodel}
Let $\fc^\circ_M\subseteq\fc_M$ be the open set where the block polynomials at
each vertex are pairwise coprime, and $\pi_M\colon\fc^\circ_M\to\fc$ the map
multiplying them.  Then $\pi_M$ is \'etale, $X(M)^\circ=\chi_M^{-1}(\fc^\circ_M)$,
and $(\iota_M,\chi_M)\colon X(M)^\circ\to X\times_\fc\fc^\circ_M$ is an
isomorphism.  Hence $\pi_M^*K\cong K_M|_{\fc^\circ_M}$.
\end{lemma}

\begin{proof}
The Jacobian determinant of $(f,g)\mapsto fg$ on pairs of monic polynomials
is, up to sign, the resultant of $f$ and $g$, so $\pi_M$ is \'etale.  The
block polynomials are coprime exactly when the block spectra are disjoint,
so $X(M)^\circ=\chi_M^{-1}(\fc^\circ_M)$.

The inverse of $(\iota_M,\chi_M)$ conjugates a point into block form along
the decomposition of each $V_k$ by the factors of $\det(t-\ell_k)$; the
off-block-diagonal part of $\ell^*$ is then the correction term of
Lemma~\ref{lem:ind}.
\begin{details}
Let $p=(\ell,\ell^*,n)$ be a
stable point of $\mu^{-1}(0)$, and let $\det(t-\ell_k)=\prod_\beta f_\beta$ be
a factorization into pairwise coprime monic polynomials, indexed by the
blocks $\beta$ at $k$, with $\deg f_\beta=|\beta|$.  By the Chinese remainder
theorem and Cayley--Hamilton, $V_k=\bigoplus_\beta\ker f_\beta(\ell_k)$, and
$\ell_k$ has characteristic polynomial $f_\beta$ on $\ker f_\beta(\ell_k)$, so
this summand has dimension $|\beta|$.  Choose $g\in G$ mapping each
$\ker f_\beta(\ell_k)$ onto the block $V^{(\beta)}$; it is unique up to left
multiplication by $M$.  Then $g\ell g^{-1}$ is block diagonal, with
disjoint block spectra.  Write $g\ell^*g^{-1}=\ell_M^*+\gamma$ with
$\ell_M^*\in\mathfrak m$ and $\gamma\in\mathfrak m^\perp$.  The
off-block-diagonal part of $\mu(gp)=0$ says $[g\ell g^{-1},\gamma]+\mathrm
pr_{\mathfrak m^\perp}\mu'(gn)=0$, so $\gamma$ is the correction term of
Lemma~\ref{lem:ind} and $gp=\tilde\iota(g\ell g^{-1},\ell_M^*,gn)$.  Its
block-diagonal part says that $(g\ell g^{-1},\ell_M^*,gn)$ lies in
$\mu_M^{-1}(0)$, and this point is stable by the stability comparison in the
proof of Lemma~\ref{lem:ind}.  Sending $([p],(f_\beta))$ to its class in
$X(M)^\circ$ is well defined, since $g$ is unique up to $M$, and it is
algebraic, since the projections onto the $\ker f_\beta(\ell_k)$ are
polynomials in $\ell_k$ with coefficients regular on $\fc_M^\circ$.  It is
inverse to $(\iota_M,\chi_M)$.
\end{details}
The last claim is base change along the \'etale map $\pi_M$.
\end{proof}

\begin{lemma}[stalks]\label{lem:stalks}
Let $M=M_m$ and $c\in\fc_m$.  Choose labeled roots $(c_D)$ of $c$, and let
$\tilde c\in\fc^\circ_M$ have block polynomials $(t-c_D)^{s_D}$, so that
$\pi_M(\tilde c)=c$.
\begin{enumerate}[leftmargin=2em]
\item There is an isomorphism $K_c\cong H^*(X(M))$ of mixed Hodge structures,
and $K_c$ is pure: $H^i(K_c)$ has weight $i$.
\item For $b\in H^*(X)$, the image of $b$ in the stalk $K_c$ vanishes if and
only if $\iota_M^*b=0$.
\end{enumerate}
\end{lemma}

\begin{proof}
Let $\mathfrak s_M\subseteq\fc_M$ be the slice where the roots of each block
polynomial sum to zero (not to be confused with the open set $\fc^\circ_M$), $X(M)_{\mathfrak s}=\chi_M^{-1}(\mathfrak s_M)$, and
$K^{\mathfrak s}_M=R\chi_{M*}\Q_{X(M)_{\mathfrak s}}$.  Translating each block of $\ell_M$ by a scalar
gives isomorphisms $X(M)\cong\mathfrak z(\mathfrak m)\times X(M)_{\mathfrak s}$ and
$\fc_M\cong\mathfrak z(\mathfrak m)\times\mathfrak s_M$ compatible with $\chi_M$, so
$K_M\cong\Q\boxtimes K^{\mathfrak s}_M$.  In these coordinates $\tilde c=((c_D),0)$.  The
complex $K^{\mathfrak s}_M$ is constructible for a stratification by
$\C^\times$-stable strata, as $\chi_M$ is $\C^\times$-equivariant, so by~\ref{I:sheaves}(d) restriction
$H^*(X(M)_{\mathfrak s})=R\Gamma(\mathfrak s_M,K^{\mathfrak s}_M)\to(K^{\mathfrak s}_M)_0$ is an isomorphism.  Since
$\pi_M$ is \'etale, Lemma~\ref{lem:localmodel} gives
\[
 K_c\cong(K_M)_{\tilde c}\cong(K^{\mathfrak s}_M)_0\cong H^*(X(M)_{\mathfrak s})\cong H^*(X(M)).
\]
These are isomorphisms of mixed Hodge structures, and $H^*(X(M))$ is pure
by~\ref{I:sheaves}(a), since $X(M)$ is smooth and retracts onto its proper
core~\ref{I:core}(a).  This proves~(1).

For~(2), let $U\ni c$ be a small neighborhood over which $\pi_M$ has a section
through $\tilde c$, with image $\tilde U$.  The image of $b$ in $K_c$ is its
restriction to $\chi^{-1}(U)$, which by Lemma~\ref{lem:localmodel} is
$\chi_M^{-1}(\tilde U)\subseteq X(M)^\circ$, embedded by $\iota_M$.  So it is
the restriction of $\iota_M^*b$ to $\chi_M^{-1}(\tilde U)$.  By
Lemma~\ref{lem:ind}(3) and the surjectivity of $\Phi$, $\iota_M^*b$ lies in
$H^*(X(M))$, and restriction from $H^*(X(M))$ to $\chi_M^{-1}(\tilde U)$ is the
isomorphism of~(1).
\end{proof}

\begin{proof}[Proof of Theorem~\ref{thm:higgs}]
\emph{Purity.}  The complex $\Q_X$ is pure of weight $0$, so $K$ has weights
$\ge0$ by~\ref{I:sheaves}(a).  By Lemma~\ref{lem:stalks}(1), every stalk
$H^i(K_c)$ has weight $i$, so $K$ has weights $\le0$.  So $K$ is pure,
and by~\ref{I:sheaves}(b) it is a direct sum of shifted simple perverse
sheaves.

\emph{Supports.}  Near a point $c\in\fc_m$, Lemma~\ref{lem:localmodel} and
the proof of Lemma~\ref{lem:stalks} identify $K$, \'etale locally, with
$\Q\boxtimes K^{\mathfrak s}_M$ near $((c_D),0)$, for $M=M_m$.  The stratum $\fc_m$
corresponds to $\mathfrak z(\mathfrak m)\times\{0\}$, since near $\tilde c$
the block polynomials stay coprime and each has a single root exactly over
$\fc_m$.  So the cohomology sheaves of $K$ are locally constant on $\fc_m$.
They become constant on the cover by labeled roots, so the monodromy factors
through $\cW_m$.
\begin{details}
That cover maps
to $\fc^\circ_M$ by $(c_D)\mapsto((t-c_D)^{s_D})_D$, which in the coordinates
$\mathfrak z(\mathfrak m)\times\mathfrak s_M$ is $(c_D)\mapsto((c_D),0)$; by
Lemma~\ref{lem:localmodel}, which holds over all of $\fc^\circ_M$, and by
$K_M\cong\Q\boxtimes K^{\mathfrak s}_M$, the pullback of $K$ along it is
constant with fiber $(K^{\mathfrak s}_M)_0$.
\end{details}
Hence each simple summand of $K$ is
$\IC(\overline{\fc_m},\cL)$ for some $m$ and some local system $\cL$ on $\fc_m$
with monodromy through $\cW_m$.  By the same argument, $K^{\mathfrak s}_M$ is pure and
is a direct sum of shifted $\IC$ sheaves with $\C^\times$-stable supports.  In
the chart, a summand of $K$ supported on $\overline{\fc_m}$
becomes $\Q$ boxed with a summand of $K^{\mathfrak s}_M$ supported at $0$, and the other
summands of $K$ through $c$ become $\Q$ boxed with summands of $K^{\mathfrak s}_M$ whose
support has positive dimension.

\emph{Multiplicities.}  Write $K^{\mathfrak s}_M=K_{\rm cusp}\oplus K'$, with $K_{\rm cusp}$ supported at $0$
and $K'$ a sum of shifted $\IC$ sheaves whose supports are
$\C^\times$-stable of positive dimension, hence contain $0$.  Restriction
from $\mathfrak s_M$ to $\mathfrak s_M\setminus0$ kills the global sections of $K_{\rm cusp}$.  It is
injective on those of $K'$, by~\ref{I:sheaves}(c) and~(d).
\begin{details}
Let $j$ be the inclusion of
$\mathfrak s_M\setminus0$ and take global sections of the triangle
$i_{0*}i_0^!\to\mathrm{id}\to j_*j^*$ applied to a summand of $K'$.  The
global sections of the summand are its stalk $i_0^*$ at $0$,
by~\ref{I:sheaves}(d), and the map from $i_0^!$ to them is the natural map
$i_0^!\to i_0^*$, which is zero by~\ref{I:sheaves}(c).  So restriction to
$\mathfrak s_M\setminus0$ is injective on that summand.
\end{details}
So the global sections of $K_{\rm cusp}$ are
the kernel of
\[
 H^*(X(M)_{\mathfrak s})\to H^*\big(X(M)_{\mathfrak s}\setminus\chi_M^{-1}(0)\big),
\]
which under $X(M)\cong\mathfrak z(\mathfrak m)\times X(M)_{\mathfrak s}$ is the kernel
$V_M$ of restriction to $X(M)\setminus\Delta_M$.  The deck transformations of
the cover by labeled roots are induced by $N_G(M)/M=\cW_m$, which permutes the
blocks, so the summand of $K$ supported on $\overline{\fc_m}$ is
$\IC(\overline{\fc_m},\cL(V_M))[-l(m)]$.  The shift $[-l(m)]$ comes from the
normalization of perverse sheaves: on $\fc_m$ this summand is $\cL(V_M)$ in the
degrees of $V_M$.  This proves~(1).

\emph{Global sections.}  The summands of $K$ are constructible for the
$\C^\times$-stable strata $\fc_m$, so by~\ref{I:sheaves}(d) their global
sections are their stalks at $0$.  The normalization
$\A^{P_m}/\cW_m\to\overline{\fc_m}$ is finite and an isomorphism over $\fc_m$, so
$\IC(\overline{\fc_m},\cL)$ is the pushforward of $\IC(\A^{P_m}/\cW_m,\cL)$, and only
$0$ lies over $0$.  So by~\ref{I:sheaves}(e), the stalk of $\IC(\overline{\fc_m},\cL(V))[-l(m)]$ at
$0$ is $V^{\cW_m}$, and it maps injectively to the stalk at every point of
$\fc_m$.  So $H^*(X)=R\Gamma(\fc,K)\cong\bigoplus_mV_{M_m}^{\cW_m}$, which is the
first claim of~(2).

Write $b\in H^*(X)$ as $\sum_mb_m$, with $b_m$ in the global sections of the
summand of $m$.  A summand of $m''$ has nonzero stalk at $c\in\fc_m$ only if
$m$ is a coarsening of $m''$.  So, by descending induction on $l(m)$, the
kernel of restriction to $\chi^{-1}\fc_{\ge L}$ is the sum of the summands
with $l(m)<L$, and it is also the set of $b$ whose stalks vanish at every
point of $\fc_{\ge L}$.
\begin{details}
The summands with $l(m)<L$ are supported on
$\overline{\fc_m}$, which consists of strata $\fc_{m'}$ with $l(m')\le l(m)$
and so does not meet $\fc_{\ge L}$; they restrict to zero on
$\chi^{-1}\fc_{\ge L}$.  Suppose $b$ restricts to zero there.  We show
$b_m=0$ for $l(m)\ge L$, by descending induction on $l(m)$.  Let $c\in\fc_m$.
A summand of $m''$ has nonzero stalk at $c$ only if $c\in\overline{\fc_{m''}}$,
that is, if $m$ is a coarsening of $m''$.  Then $l(m'')\ge l(m)$, with equality
only for $m''=m$.  By induction $b_{m''}=0$ when $l(m'')>l(m)$, so the stalk of
$b$ at $c$ is that of $b_m$.  It vanishes, so $b_m=0$ by the injectivity
above.  Hence the kernel of restriction to $\chi^{-1}\fc_{\ge L}$ is the sum
of the summands with $l(m)<L$, and it is also the set of $b$ whose stalks
vanish at every point of $\fc_{\ge L}$.
\end{details}
By Lemma~\ref{lem:stalks}(2), the
stalks of $b$ over $\fc_{m'}$ vanish if and only if $\iota_{M_{m'}}^*b=0$.
Every $M\in\mathcal M_L$ is conjugate to some $M_{m'}$ with $l(m')\ge L$, so
this kernel is $\bigcap_{M\in\mathcal M_L}\ker\iota_M^*$.  Finally, by
Lemma~\ref{lem:ind}(3), $\iota_M^*\circ\Phi=\Phi_M\circ\res_M$, and $\Phi_M$ is
injective by Theorem~\ref{thm:loops} for $Q_M$.  So
$\bigcap_{M\in\mathcal M_L}\ker\iota_M^*=\Phi\big(\bigcap_{M\in\mathcal
M_L}\ker\res_M\big)=\Phi(F_{L-1})$ by Corollary~\ref{cor:cusp}.  This
proves~(2).

\emph{Cuspidal parts.}  Apply~(2) to $X(M)$ for $M=M_m$, with $l(m)+1$ in
place of $L$.  The vertices of $Q_M$ are the $l(m)$ blocks, so the loops of
$Q_M$ have at least $l(m)+1$ distinct eigenvalues in all exactly off
$\Delta_M$.  So $V_M=\Phi_M(F_{l(m)}(M))$.  By Corollary~\ref{cor:cusp} for
$Q_M$, whose gauge group is $M$, with $l(m)$ blocks, $F_{l(m)}(M)$ is the
cuspidal part $(\Pol_m/J_m)\langle\delta_m\rangle$ of $B(M)$.  The group $N_G(M)/M=\cW_m$ permutes the blocks of
equal size at each vertex, so it acts on the gauge theory of $Q_M$, hence on
$B(M)$, $H^*(X(M))$ and $\Pol_m$, and $\Phi_M$ and $\epsilon$ commute with these
actions.  So the identification is $\cW_m$-equivariant.  This proves~(3).
\end{proof}

\begin{remark}
For the Jordan quiver with $\bw=1$, $\chi$ is the composite
$\Hilb^n(\C^2)\to\mathrm{Sym}^n\C^2\to\mathrm{Sym}^n\C$, every $V_{M_m}$ is
one-dimensional, and Theorem~\ref{thm:higgs}(1) is the push-forward of the
decomposition of G\"ottsche and Soergel for the Hilbert--Chow map.
\end{remark}

\begin{remark}
Decompositions over multiplicity strata also appear for the BPS sheaves of
preprojective algebras, which for the Jordan quiver are supported on small
diagonals \cite[Thm.~5.1, Cor.~5.3]{davisonIntegralityConjecture2023}.  For any quiver,
\cite[Thm.~1.14]{davisonBPSAlgebras2023} decomposes the cohomology of
Nakajima quiver varieties over the BPS Lie algebra; that decomposition is along
$X\to\mu^{-1}(0)/\!/G$ rather than along $\chi$.  We have not compared these
with Theorem~\ref{thm:higgs}.
\end{remark}

\section{Status}\label{sec:status}
\begin{enumerate}[leftmargin=2em]
\item Theorem~\ref{thm:map} is announced in \cite[Thm.~E]{bottaQuasimapCritical2026}; we do
not give a proof.
\item We have not checked the following citations against the sources:
\begin{itemize}[leftmargin=1.5em]
\item the form of \cite[Prop.~3.22]{BFN} used in~\ref{I:base}, with
flavor and loop rotation;
\item \cite{hauselCohomologyLarge2013} and \cite{greeneInterpretationWhitney1983};
\item \cite[Prop.~3.18]{bellamySymplecticResolutions2021} in~\ref{I:rot};
\item the theorem numbers in the paragraph on known cases (\S\ref{ss:setup}) and
in the remarks after Theorem~\ref{thm:higgs}, which were taken from the arXiv
versions;
\item the smoothness statement in~\ref{I:hypertoric}(a), for which we give
no theorem number;
\item the degree of $r_\lambda$ in the grading of \cite{BFN}
(\S\ref{ss:abel});
\item the form with flavor of \cite[Thm.~4.1, Lemma~5.17,
Prop.~6.6]{BFN}, used before Proposition~\ref{prop:loopremoval}; these
were checked without flavor against arXiv:1601.03586v8;
\item \cite[\S2.6]{CG97} for the top homology of a compact variety
in~\ref{I:core}(d);
\item \cite[Prop.~3.7.5]{kashiwaraSheavesManifolds1994};
\item the numbering in \cite{BBD};
\item that \cite[Thm.~4.1(2)]{nakajimaInstantonsALE1994}, whose proof there
refers to Kronheimer, holds for quivers with loops (remark after
Lemma~\ref{lem:betti}).
\end{itemize}
\item The second proof in \S\ref{sec:bps} depends on~\ref{H:poly}
and~\ref{H:bps}, whose proofs in \cite{jindalLoopNilpotent2026} are only
sketched (Remark~\ref{rem:bps-status}), and on \pre{Lemma~8.1}.  The sketch
of Lemma~\ref{lem:hall-coulomb} follows \pre{Lemma~8.1}.  Its signs and
specialization were compared with \cite{tsymbaliukDifferenceOperators2023}
only in a machine reading of the preprint, and the formula \pre{(8.2)} was
not rederived here.
\item A proof on the quiver variety side that the cuspidal part $V_G$ is
$(\Pol_G/J_G)\langle\delta_{m_G}\rangle$ would give the formula of
Theorem~\ref{thm:loops} with no input from the Coulomb branch.  Hausel's
formula and Remark~\ref{rem:kirwan}(4) already give that formula with no such
input (Remark~\ref{rem:hausel}), and Theorem~\ref{thm:higgs}(2) gives
$H^*(X)\cong\bigoplus_mV_{M_m}^{\cW_m}$ without using $\Phi$; what is missing is
the comparison summand by summand.
\end{enumerate}

% \markboth is neutralized so biblatex does not overwrite the running head.
{\renewcommand{\markboth}[2]{}\printbibliography}

\end{document}
